\documentclass[11pt]{article}
\usepackage[margin=1in]{geometry}
\usepackage{amsmath,amssymb,amsthm}
\usepackage{booktabs}
\usepackage{hyperref}
% keep build dates, time zone and included-file paths out of the PDF (anonymous review)
\pdfinfoomitdate=1
\pdfsuppressptexinfo=-1
\usepackage{enumitem}
% generated by artifact/src/render.py, do not edit
\newcommand{\EOneAudCIhi}{8.1}
\newcommand{\EOneAudCIhiZero}{1.6}
\newcommand{\EOneAudCIlo}{8.1}
\newcommand{\EOneAudCIloZero}{1.5}
\newcommand{\EOneAudHalfOne}{3.8}
\newcommand{\EOneAudHalfTenth}{4.5}
\newcommand{\EOneAudHalfZero}{2.3}
\newcommand{\EOneAudPhiOneRTen}{9.5}
\newcommand{\EOneAudPhiOneRZero}{5.3}
\newcommand{\EOneAudPhiZeroRTen}{8.1}
\newcommand{\EOneAudPhiZeroRZero}{1.5}
\newcommand{\EOneAudRatioPhiOne}{1.8}
\newcommand{\EOneAudRatioPhiZero}{5.2}
\newcommand{\EOneBlindRecHalfTenth}{0.04}
\newcommand{\EOneBlindRecHalfZero}{0.35}
\newcommand{\EOneBlindRecMax}{0.35}
\newcommand{\EOneDPIrr}{3.08}
\newcommand{\EOneEpisodes}{50{,}000}
\newcommand{\EOneGapBase}{6.5}
\newcommand{\EOneGapMax}{69.6}
\newcommand{\EOneGapMaxPhi}{0.0}
\newcommand{\EOneGapMaxRone}{0}
\newcommand{\EOneGrid}{1{,}001}
\newcommand{\EOneIrrBlindIrr}{6.98}
\newcommand{\EOneIrrBlindRegretMax}{38.7}
\newcommand{\EOneIrrBlindRegretTen}{0.3}
\newcommand{\EOneLamRegretMax}{3.5}
\newcommand{\EOneMyDprecTen}{5.1}
\newcommand{\EOneMyTen}{55.3}
\newcommand{\EOneNoAuditTen}{53.4}
\newcommand{\EOneRegretMin}{-1.1}
\newcommand{\EOneRollbackBlindMax}{16.7}
\newcommand{\EOneRollbackBlindMin}{4.5}
\newcommand{\EOneRollbackOptimistMax}{12.8}
\newcommand{\EOneSeeds}{5}
\newcommand{\EOneSensBreaks}{$c_a=2$, $d=0.5$, $d=0.7$ and $\lambda=0.02$}
\newcommand{\EOneSensDhighAud}{7.6}
\newcommand{\EOneSensDlowAud}{3.8}
\newcommand{\EOneSensGapMax}{237.9}
\newcommand{\EOneSensGapMin}{0.0}
\newcommand{\EOneSensHolds}{9}
\newcommand{\EOneSensShighAud}{4.1}
\newcommand{\EOneSensSlowAud}{10.5}
\newcommand{\EOneSensTotal}{13}
\newcommand{\EOneSensorRegretMax}{5.8}
\newcommand{\EOneStopRule}{passes}
\newcommand{\EThreeAdBase}{immutable validator}
\newcommand{\EThreeAdBaseC}{memory}
\newcommand{\EThreeAdBaseHi}{11.7}
\newcommand{\EThreeAdBaseLo}{8.6}
\newcommand{\EThreeAdBaseV}{10.0}
\newcommand{\EThreeAdCo}{11.8}
\newcommand{\EThreeAdCoC}{memory note}
\newcommand{\EThreeAdNo}{14.1}
\newcommand{\EThreeAdNoC}{skill}
\newcommand{\EThreeAdRb}{14.0}
\newcommand{\EThreeAdRbC}{memory note}
\newcommand{\EThreeAdRover}{11.2}
\newcommand{\EThreeAdRoverC}{memory}
\newcommand{\EThreeAdRoverCN}{memory note}
\newcommand{\EThreeAdRoverHi}{12.6}
\newcommand{\EThreeAdRoverLo}{10.0}
\newcommand{\EThreeAllocBranchTen}{0.29}
\newcommand{\EThreeAllocBranchZero}{0.04}
\newcommand{\EThreeAllocMemoryTen}{0.17}
\newcommand{\EThreeAllocMemoryZero}{0.05}
\newcommand{\EThreeAllocQueueTen}{0.24}
\newcommand{\EThreeAllocQueueZero}{0.14}
\newcommand{\EThreeAllocSkillTen}{0.30}
\newcommand{\EThreeAllocSkillZero}{0.28}
\newcommand{\EThreeAllocTotTen}{1.00}
\newcommand{\EThreeAllocTotZero}{0.51}
\newcommand{\EThreeBestBase}{rollback only}
\newcommand{\EThreeBestBaseHi}{13.09}
\newcommand{\EThreeBestBaseLo}{9.75}
\newcommand{\EThreeBestBaseLoss}{11.36}
\newcommand{\EThreeCBlAud}{-0.62}
\newcommand{\EThreeCBlAudHi}{-0.53}
\newcommand{\EThreeCBlAudLo}{-0.72}
\newcommand{\EThreeCBlAudP}{<0.001}
\newcommand{\EThreeCBlEv}{0.11}
\newcommand{\EThreeCBlEvHi}{0.15}
\newcommand{\EThreeCBlEvLo}{0.07}
\newcommand{\EThreeCBlEvP}{<0.001}
\newcommand{\EThreeCBlFr}{-0.03}
\newcommand{\EThreeCBlFrHi}{-0.01}
\newcommand{\EThreeCBlFrLo}{-0.05}
\newcommand{\EThreeCBlFrP}{=0.008}
\newcommand{\EThreeCBlIrr}{2.17}
\newcommand{\EThreeCBlIrrHi}{2.92}
\newcommand{\EThreeCBlIrrLo}{1.50}
\newcommand{\EThreeCBlIrrP}{<0.001}
\newcommand{\EThreeCBlLoss}{1.90}
\newcommand{\EThreeCBlLossHi}{2.82}
\newcommand{\EThreeCBlLossLo}{1.05}
\newcommand{\EThreeCBlLossP}{<0.001}
\newcommand{\EThreeCBlRep}{0.02}
\newcommand{\EThreeCBlRepHi}{0.48}
\newcommand{\EThreeCBlRepLo}{-0.49}
\newcommand{\EThreeCBlRepP}{=0.745}
\newcommand{\EThreeCBlRes}{0.34}
\newcommand{\EThreeCBlResHi}{0.60}
\newcommand{\EThreeCBlResLo}{0.12}
\newcommand{\EThreeCBlResP}{=0.005}
\newcommand{\EThreeCCoAud}{0.05}
\newcommand{\EThreeCCoAudHi}{0.10}
\newcommand{\EThreeCCoAudLo}{-0.00}
\newcommand{\EThreeCCoAudP}{=0.094}
\newcommand{\EThreeCCoEv}{0.01}
\newcommand{\EThreeCCoEvHi}{0.02}
\newcommand{\EThreeCCoEvLo}{0.00}
\newcommand{\EThreeCCoEvP}{=0.157}
\newcommand{\EThreeCCoFr}{0.01}
\newcommand{\EThreeCCoFrHi}{0.03}
\newcommand{\EThreeCCoFrLo}{-0.00}
\newcommand{\EThreeCCoFrP}{=0.180}
\newcommand{\EThreeCCoIrr}{0.17}
\newcommand{\EThreeCCoIrrHi}{0.42}
\newcommand{\EThreeCCoIrrLo}{0.00}
\newcommand{\EThreeCCoIrrP}{=0.157}
\newcommand{\EThreeCCoLoss}{0.58}
\newcommand{\EThreeCCoLossHi}{1.01}
\newcommand{\EThreeCCoLossLo}{0.25}
\newcommand{\EThreeCCoLossP}{<0.001}
\newcommand{\EThreeCCoRep}{0.29}
\newcommand{\EThreeCCoRepHi}{0.65}
\newcommand{\EThreeCCoRepLo}{-0.05}
\newcommand{\EThreeCCoRepP}{=0.018}
\newcommand{\EThreeCCoRes}{0.08}
\newcommand{\EThreeCCoResHi}{0.24}
\newcommand{\EThreeCCoResLo}{0.00}
\newcommand{\EThreeCCoResP}{=0.317}
\newcommand{\EThreeCMyAud}{-0.77}
\newcommand{\EThreeCMyAudHi}{-0.66}
\newcommand{\EThreeCMyAudLo}{-0.88}
\newcommand{\EThreeCMyAudP}{<0.001}
\newcommand{\EThreeCMyEv}{0.10}
\newcommand{\EThreeCMyEvHi}{0.14}
\newcommand{\EThreeCMyEvLo}{0.07}
\newcommand{\EThreeCMyEvP}{<0.001}
\newcommand{\EThreeCMyFr}{-0.03}
\newcommand{\EThreeCMyFrHi}{-0.01}
\newcommand{\EThreeCMyFrLo}{-0.06}
\newcommand{\EThreeCMyFrP}{=0.011}
\newcommand{\EThreeCMyIrr}{2.00}
\newcommand{\EThreeCMyIrrHi}{2.75}
\newcommand{\EThreeCMyIrrLo}{1.33}
\newcommand{\EThreeCMyIrrP}{<0.001}
\newcommand{\EThreeCMyLoss}{2.22}
\newcommand{\EThreeCMyLossHi}{3.26}
\newcommand{\EThreeCMyLossLo}{1.21}
\newcommand{\EThreeCMyLossP}{<0.001}
\newcommand{\EThreeCMyRep}{0.36}
\newcommand{\EThreeCMyRepHi}{1.14}
\newcommand{\EThreeCMyRepLo}{-0.45}
\newcommand{\EThreeCMyRepP}{=0.475}
\newcommand{\EThreeCMyRes}{0.62}
\newcommand{\EThreeCMyResHi}{1.07}
\newcommand{\EThreeCMyResLo}{0.23}
\newcommand{\EThreeCMyResP}{=0.005}
\newcommand{\EThreeCNaAud}{-0.90}
\newcommand{\EThreeCNaAudHi}{-0.78}
\newcommand{\EThreeCNaAudLo}{-1.02}
\newcommand{\EThreeCNaAudP}{<0.001}
\newcommand{\EThreeCNaEv}{0.05}
\newcommand{\EThreeCNaEvH}{5.4}
\newcommand{\EThreeCNaEvHHi}{8.3}
\newcommand{\EThreeCNaEvHLo}{2.9}
\newcommand{\EThreeCNaEvHi}{0.08}
\newcommand{\EThreeCNaEvLo}{0.03}
\newcommand{\EThreeCNaEvOther}{0}
\newcommand{\EThreeCNaEvP}{=0.002}
\newcommand{\EThreeCNaFr}{0.03}
\newcommand{\EThreeCNaFrHi}{0.07}
\newcommand{\EThreeCNaFrLo}{-0.01}
\newcommand{\EThreeCNaFrP}{=0.210}
\newcommand{\EThreeCNaIrr}{1.08}
\newcommand{\EThreeCNaIrrHi}{1.67}
\newcommand{\EThreeCNaIrrLo}{0.58}
\newcommand{\EThreeCNaIrrP}{=0.002}
\newcommand{\EThreeCNaLoss}{1.09}
\newcommand{\EThreeCNaLossHi}{1.82}
\newcommand{\EThreeCNaLossLo}{0.39}
\newcommand{\EThreeCNaLossP}{=0.008}
\newcommand{\EThreeCNaRep}{0.55}
\newcommand{\EThreeCNaRepHi}{1.08}
\newcommand{\EThreeCNaRepLo}{-0.00}
\newcommand{\EThreeCNaRepP}{=0.060}
\newcommand{\EThreeCNaRes}{0.35}
\newcommand{\EThreeCNaResHi}{0.65}
\newcommand{\EThreeCNaResLo}{0.12}
\newcommand{\EThreeCNaResP}{=0.012}
\newcommand{\EThreeCNoAud}{-0.90}
\newcommand{\EThreeCNoAudHi}{-0.78}
\newcommand{\EThreeCNoAudLo}{-1.02}
\newcommand{\EThreeCNoAudP}{<0.001}
\newcommand{\EThreeCNoEv}{0.13}
\newcommand{\EThreeCNoEvHi}{0.17}
\newcommand{\EThreeCNoEvLo}{0.10}
\newcommand{\EThreeCNoEvP}{<0.001}
\newcommand{\EThreeCNoFr}{-0.03}
\newcommand{\EThreeCNoFrHi}{-0.01}
\newcommand{\EThreeCNoFrLo}{-0.06}
\newcommand{\EThreeCNoFrP}{=0.011}
\newcommand{\EThreeCNoIrr}{2.67}
\newcommand{\EThreeCNoIrrHi}{3.50}
\newcommand{\EThreeCNoIrrLo}{1.92}
\newcommand{\EThreeCNoIrrP}{<0.001}
\newcommand{\EThreeCNoLoss}{2.58}
\newcommand{\EThreeCNoLossHi}{3.71}
\newcommand{\EThreeCNoLossLo}{1.57}
\newcommand{\EThreeCNoLossP}{<0.001}
\newcommand{\EThreeCNoRep}{-0.27}
\newcommand{\EThreeCNoRepHi}{0.44}
\newcommand{\EThreeCNoRepLo}{-1.02}
\newcommand{\EThreeCNoRepP}{=0.552}
\newcommand{\EThreeCNoRes}{1.08}
\newcommand{\EThreeCNoResHi}{1.56}
\newcommand{\EThreeCNoResLo}{0.66}
\newcommand{\EThreeCNoResP}{<0.001}
\newcommand{\EThreeCRbAud}{-0.90}
\newcommand{\EThreeCRbAudHi}{-0.78}
\newcommand{\EThreeCRbAudLo}{-1.02}
\newcommand{\EThreeCRbAudP}{<0.001}
\newcommand{\EThreeCRbEv}{0.05}
\newcommand{\EThreeCRbEvH}{4.6}
\newcommand{\EThreeCRbEvHHi}{7.1}
\newcommand{\EThreeCRbEvHLo}{2.5}
\newcommand{\EThreeCRbEvHi}{0.07}
\newcommand{\EThreeCRbEvLo}{0.03}
\newcommand{\EThreeCRbEvOther}{0}
\newcommand{\EThreeCRbEvP}{=0.002}
\newcommand{\EThreeCRbFr}{0.05}
\newcommand{\EThreeCRbFrHi}{0.10}
\newcommand{\EThreeCRbFrLo}{0.00}
\newcommand{\EThreeCRbFrP}{=0.042}
\newcommand{\EThreeCRbIrr}{0.92}
\newcommand{\EThreeCRbIrrHi}{1.42}
\newcommand{\EThreeCRbIrrLo}{0.50}
\newcommand{\EThreeCRbIrrP}{=0.002}
\newcommand{\EThreeCRbLoss}{2.08}
\newcommand{\EThreeCRbLossHi}{2.93}
\newcommand{\EThreeCRbLossLo}{1.28}
\newcommand{\EThreeCRbLossP}{<0.001}
\newcommand{\EThreeCRbRep}{2.31}
\newcommand{\EThreeCRbRepHi}{3.01}
\newcommand{\EThreeCRbRepLo}{1.60}
\newcommand{\EThreeCRbRepP}{<0.001}
\newcommand{\EThreeCRbRes}{-0.25}
\newcommand{\EThreeCRbResHi}{0.01}
\newcommand{\EThreeCRbResLo}{-0.51}
\newcommand{\EThreeCRbResP}{=0.092}
\newcommand{\EThreeCanaryRoverMain}{0.17}
\newcommand{\EThreeCherry}{22{,}569}
\newcommand{\EThreeCommits}{328{,}320}
\newcommand{\EThreeConflicts}{4{,}448}
\newcommand{\EThreeDOneBlind}{1.90}
\newcommand{\EThreeDOneBlindHi}{2.88}
\newcommand{\EThreeDOneBlindLo}{1.04}
\newcommand{\EThreeDOneConst}{0.58}
\newcommand{\EThreeDOneConstHi}{0.98}
\newcommand{\EThreeDOneConstLo}{0.23}
\newcommand{\EThreeDOneMyopic}{2.22}
\newcommand{\EThreeDOneMyopicHi}{3.31}
\newcommand{\EThreeDOneMyopicLo}{1.24}
\newcommand{\EThreeDOneNoAudit}{2.58}
\newcommand{\EThreeDOneNoAuditHi}{3.80}
\newcommand{\EThreeDOneNoAuditLo}{1.60}
\newcommand{\EThreeDOneNoaudit}{1.09}
\newcommand{\EThreeDOneNoauditHi}{1.85}
\newcommand{\EThreeDOneNoauditLo}{0.39}
\newcommand{\EThreeDOnePeriodic}{6.79}
\newcommand{\EThreeDOnePeriodicHi}{7.75}
\newcommand{\EThreeDOnePeriodicLo}{5.90}
\newcommand{\EThreeDOneRandom}{4.18}
\newcommand{\EThreeDOneRandomHi}{5.27}
\newcommand{\EThreeDOneRandomLo}{3.18}
\newcommand{\EThreeDOneReset}{16.66}
\newcommand{\EThreeDOneResetHi}{18.59}
\newcommand{\EThreeDOneResetLo}{14.70}
\newcommand{\EThreeDOneRollback}{2.08}
\newcommand{\EThreeDOneRollbackHi}{2.95}
\newcommand{\EThreeDOneRollbackLo}{1.27}
\newcommand{\EThreeDOneValidator}{5.91}
\newcommand{\EThreeDOneValidatorHi}{6.86}
\newcommand{\EThreeDOneValidatorLo}{4.94}
\newcommand{\EThreeDTenBlind}{1.17}
\newcommand{\EThreeDTenBlindHi}{2.49}
\newcommand{\EThreeDTenBlindLo}{-0.05}
\newcommand{\EThreeDTenConst}{-0.03}
\newcommand{\EThreeDTenConstHi}{0.57}
\newcommand{\EThreeDTenConstLo}{-0.65}
\newcommand{\EThreeDTenMyopic}{3.55}
\newcommand{\EThreeDTenMyopicHi}{5.10}
\newcommand{\EThreeDTenMyopicLo}{2.04}
\newcommand{\EThreeDTenNoAudit}{1.87}
\newcommand{\EThreeDTenNoAuditHi}{3.62}
\newcommand{\EThreeDTenNoAuditLo}{0.19}
\newcommand{\EThreeDTenNoaudit}{1.12}
\newcommand{\EThreeDTenNoauditHi}{2.62}
\newcommand{\EThreeDTenNoauditLo}{-0.35}
\newcommand{\EThreeDTenPeriodic}{7.23}
\newcommand{\EThreeDTenPeriodicHi}{8.87}
\newcommand{\EThreeDTenPeriodicLo}{5.51}
\newcommand{\EThreeDTenRandom}{4.54}
\newcommand{\EThreeDTenRandomHi}{6.52}
\newcommand{\EThreeDTenRandomLo}{2.54}
\newcommand{\EThreeDTenReplayOnly}{20.82}
\newcommand{\EThreeDTenReplayOnlyHi}{26.43}
\newcommand{\EThreeDTenReplayOnlyLo}{16.00}
\newcommand{\EThreeDTenReset}{7.11}
\newcommand{\EThreeDTenResetHi}{8.98}
\newcommand{\EThreeDTenResetLo}{5.42}
\newcommand{\EThreeDTenRollback}{4.20}
\newcommand{\EThreeDTenRollbackHi}{5.83}
\newcommand{\EThreeDTenRollbackLo}{2.68}
\newcommand{\EThreeDTenValidator}{3.95}
\newcommand{\EThreeDTenValidatorHi}{5.57}
\newcommand{\EThreeDTenValidatorLo}{2.29}
\newcommand{\EThreeDTenthBlind}{1.22}
\newcommand{\EThreeDTenthBlindHi}{1.99}
\newcommand{\EThreeDTenthBlindLo}{0.51}
\newcommand{\EThreeDTenthConst}{0.77}
\newcommand{\EThreeDTenthConstHi}{1.05}
\newcommand{\EThreeDTenthConstLo}{0.47}
\newcommand{\EThreeDTenthMyopic}{2.35}
\newcommand{\EThreeDTenthMyopicHi}{3.29}
\newcommand{\EThreeDTenthMyopicLo}{1.48}
\newcommand{\EThreeDTenthNoAudit}{2.70}
\newcommand{\EThreeDTenthNoAuditHi}{3.60}
\newcommand{\EThreeDTenthNoAuditLo}{1.84}
\newcommand{\EThreeDTenthNoaudit}{0.69}
\newcommand{\EThreeDTenthNoauditHi}{1.19}
\newcommand{\EThreeDTenthNoauditLo}{0.18}
\newcommand{\EThreeDTenthPeriodic}{6.66}
\newcommand{\EThreeDTenthPeriodicHi}{7.39}
\newcommand{\EThreeDTenthPeriodicLo}{5.93}
\newcommand{\EThreeDTenthRandom}{3.32}
\newcommand{\EThreeDTenthRandomHi}{4.13}
\newcommand{\EThreeDTenthRandomLo}{2.59}
\newcommand{\EThreeDTenthReset}{20.54}
\newcommand{\EThreeDTenthResetHi}{22.81}
\newcommand{\EThreeDTenthResetLo}{18.30}
\newcommand{\EThreeDTenthRollback}{0.53}
\newcommand{\EThreeDTenthRollbackHi}{1.19}
\newcommand{\EThreeDTenthRollbackLo}{-0.07}
\newcommand{\EThreeDTenthValidator}{5.78}
\newcommand{\EThreeDTenthValidatorHi}{6.50}
\newcommand{\EThreeDTenthValidatorLo}{5.06}
\newcommand{\EThreeDZeroBlind}{1.64}
\newcommand{\EThreeDZeroBlindHi}{2.42}
\newcommand{\EThreeDZeroBlindLo}{0.92}
\newcommand{\EThreeDZeroConst}{1.05}
\newcommand{\EThreeDZeroConstHi}{1.46}
\newcommand{\EThreeDZeroConstLo}{0.60}
\newcommand{\EThreeDZeroMyopic}{2.49}
\newcommand{\EThreeDZeroMyopicHi}{3.47}
\newcommand{\EThreeDZeroMyopicLo}{1.57}
\newcommand{\EThreeDZeroNoAudit}{2.83}
\newcommand{\EThreeDZeroNoAuditHi}{3.78}
\newcommand{\EThreeDZeroNoAuditLo}{1.94}
\newcommand{\EThreeDZeroNoaudit}{0.56}
\newcommand{\EThreeDZeroNoauditHi}{1.07}
\newcommand{\EThreeDZeroNoauditLo}{0.06}
\newcommand{\EThreeDZeroPeriodic}{6.76}
\newcommand{\EThreeDZeroPeriodicHi}{7.53}
\newcommand{\EThreeDZeroPeriodicLo}{5.98}
\newcommand{\EThreeDZeroRandom}{3.29}
\newcommand{\EThreeDZeroRandomHi}{4.15}
\newcommand{\EThreeDZeroRandomLo}{2.48}
\newcommand{\EThreeDZeroReset}{21.09}
\newcommand{\EThreeDZeroResetHi}{23.43}
\newcommand{\EThreeDZeroResetLo}{18.82}
\newcommand{\EThreeDZeroRollback}{0.48}
\newcommand{\EThreeDZeroRollbackHi}{1.19}
\newcommand{\EThreeDZeroRollbackLo}{-0.17}
\newcommand{\EThreeDZeroValidator}{5.88}
\newcommand{\EThreeDZeroValidatorHi}{6.61}
\newcommand{\EThreeDZeroValidatorLo}{5.13}
\newcommand{\EThreeEpisodes}{13{,}680}
\newcommand{\EThreeEpisodesAll}{29{,}520}
\newcommand{\EThreeExtraEpisodes}{8{,}160}
\newcommand{\EThreeFamMin}{flat}
\newcommand{\EThreeFamMinD}{0.95}
\newcommand{\EThreeFamMinHi}{2.40}
\newcommand{\EThreeFamMinLo}{-0.09}
\newcommand{\EThreeFreePe}{10.07}
\newcommand{\EThreeFreePeD}{0.79}
\newcommand{\EThreeFreePeHi}{1.74}
\newcommand{\EThreeFreePeLo}{-0.05}
\newcommand{\EThreeFreeVa}{9.64}
\newcommand{\EThreeFreeVaD}{0.36}
\newcommand{\EThreeFreeVaHi}{1.31}
\newcommand{\EThreeFreeVaLo}{-0.51}
\newcommand{\EThreeGPAwareMax}{0.15}
\newcommand{\EThreeGPAwareMin}{0.14}
\newcommand{\EThreeGPFreqFullAFifth}{9.23}
\newcommand{\EThreeGPFreqFullAFifthD}{-0.04}
\newcommand{\EThreeGPFreqFullAFifthHi}{0.77}
\newcommand{\EThreeGPFreqFullAFifthLo}{-0.92}
\newcommand{\EThreeGPFreqFullAHalf}{11.07}
\newcommand{\EThreeGPFreqFullAHalfD}{1.80}
\newcommand{\EThreeGPFreqFullAHalfHi}{2.63}
\newcommand{\EThreeGPFreqFullAHalfLo}{0.91}
\newcommand{\EThreeGPFreqFullAOne}{14.14}
\newcommand{\EThreeGPFreqFullAOneD}{4.86}
\newcommand{\EThreeGPFreqFullAOneHi}{5.74}
\newcommand{\EThreeGPFreqFullAOneLo}{3.95}
\newcommand{\EThreeGPFreqFullATenth}{8.62}
\newcommand{\EThreeGPFreqFullATenthD}{-0.66}
\newcommand{\EThreeGPFreqFullATenthHi}{0.16}
\newcommand{\EThreeGPFreqFullATenthLo}{-1.53}
\newcommand{\EThreeGPFreqFullAmGFifth}{-0.21}
\newcommand{\EThreeGPFreqFullAmGFifthHi}{0.10}
\newcommand{\EThreeGPFreqFullAmGFifthLo}{-0.56}
\newcommand{\EThreeGPFreqFullAmGHalf}{-0.22}
\newcommand{\EThreeGPFreqFullAmGHalfHi}{0.09}
\newcommand{\EThreeGPFreqFullAmGHalfLo}{-0.57}
\newcommand{\EThreeGPFreqFullAmGOne}{-0.24}
\newcommand{\EThreeGPFreqFullAmGOneHi}{0.07}
\newcommand{\EThreeGPFreqFullAmGOneLo}{-0.60}
\newcommand{\EThreeGPFreqFullAmGTenth}{-0.21}
\newcommand{\EThreeGPFreqFullAmGTenthHi}{0.10}
\newcommand{\EThreeGPFreqFullAmGTenthLo}{-0.55}
\newcommand{\EThreeGPFreqFullBreakeven}{0.17}
\newcommand{\EThreeGPFreqFullCanaryAware}{4.55}
\newcommand{\EThreeGPFreqFullCanaryGate}{7.41}
\newcommand{\EThreeGPFreqFullCanaryRG}{2.94}
\newcommand{\EThreeGPFreqFullChecks}{6.2}
\newcommand{\EThreeGPFreqFullGFifth}{9.45}
\newcommand{\EThreeGPFreqFullGFifthD}{0.17}
\newcommand{\EThreeGPFreqFullGFifthHi}{1.18}
\newcommand{\EThreeGPFreqFullGFifthLo}{-0.82}
\newcommand{\EThreeGPFreqFullGHalf}{11.30}
\newcommand{\EThreeGPFreqFullGHalfD}{2.02}
\newcommand{\EThreeGPFreqFullGHalfHi}{3.06}
\newcommand{\EThreeGPFreqFullGHalfLo}{1.02}
\newcommand{\EThreeGPFreqFullGOne}{14.38}
\newcommand{\EThreeGPFreqFullGOneD}{5.10}
\newcommand{\EThreeGPFreqFullGOneHi}{6.17}
\newcommand{\EThreeGPFreqFullGOneLo}{4.08}
\newcommand{\EThreeGPFreqFullGTenth}{8.83}
\newcommand{\EThreeGPFreqFullGTenthD}{-0.45}
\newcommand{\EThreeGPFreqFullGTenthHi}{0.55}
\newcommand{\EThreeGPFreqFullGTenthLo}{-1.43}
\newcommand{\EThreeGPFreqFullUFifth}{9.17}
\newcommand{\EThreeGPFreqFullUFifthD}{-0.11}
\newcommand{\EThreeGPFreqFullUFifthHi}{0.60}
\newcommand{\EThreeGPFreqFullUFifthLo}{-0.89}
\newcommand{\EThreeGPFreqFullUHalf}{10.98}
\newcommand{\EThreeGPFreqFullUHalfD}{1.70}
\newcommand{\EThreeGPFreqFullUHalfHi}{2.43}
\newcommand{\EThreeGPFreqFullUHalfLo}{0.93}
\newcommand{\EThreeGPFreqFullUOne}{14.00}
\newcommand{\EThreeGPFreqFullUOneD}{4.72}
\newcommand{\EThreeGPFreqFullUOneHi}{5.48}
\newcommand{\EThreeGPFreqFullUOneLo}{3.95}
\newcommand{\EThreeGPFreqFullUTenth}{8.57}
\newcommand{\EThreeGPFreqFullUTenthD}{-0.71}
\newcommand{\EThreeGPFreqFullUTenthHi}{-0.00}
\newcommand{\EThreeGPFreqFullUTenthLo}{-1.50}
\newcommand{\EThreeGPFreqHalfAFifth}{10.30}
\newcommand{\EThreeGPFreqHalfAFifthD}{1.02}
\newcommand{\EThreeGPFreqHalfAFifthHi}{1.67}
\newcommand{\EThreeGPFreqHalfAFifthLo}{0.35}
\newcommand{\EThreeGPFreqHalfAHalf}{12.10}
\newcommand{\EThreeGPFreqHalfAHalfD}{2.82}
\newcommand{\EThreeGPFreqHalfAHalfHi}{3.51}
\newcommand{\EThreeGPFreqHalfAHalfLo}{2.12}
\newcommand{\EThreeGPFreqHalfAOne}{15.10}
\newcommand{\EThreeGPFreqHalfAOneD}{5.82}
\newcommand{\EThreeGPFreqHalfAOneHi}{6.61}
\newcommand{\EThreeGPFreqHalfAOneLo}{5.08}
\newcommand{\EThreeGPFreqHalfATenth}{9.70}
\newcommand{\EThreeGPFreqHalfATenthD}{0.42}
\newcommand{\EThreeGPFreqHalfATenthHi}{1.06}
\newcommand{\EThreeGPFreqHalfATenthLo}{-0.24}
\newcommand{\EThreeGPFreqHalfAmGFifth}{-1.29}
\newcommand{\EThreeGPFreqHalfAmGFifthHi}{-0.53}
\newcommand{\EThreeGPFreqHalfAmGFifthLo}{-2.15}
\newcommand{\EThreeGPFreqHalfAmGHalf}{-1.30}
\newcommand{\EThreeGPFreqHalfAmGHalfHi}{-0.54}
\newcommand{\EThreeGPFreqHalfAmGHalfLo}{-2.17}
\newcommand{\EThreeGPFreqHalfAmGOne}{-1.33}
\newcommand{\EThreeGPFreqHalfAmGOneHi}{-0.57}
\newcommand{\EThreeGPFreqHalfAmGOneLo}{-2.20}
\newcommand{\EThreeGPFreqHalfAmGTenth}{-1.28}
\newcommand{\EThreeGPFreqHalfAmGTenthHi}{-0.52}
\newcommand{\EThreeGPFreqHalfAmGTenthLo}{-2.14}
\newcommand{\EThreeGPFreqHalfBreakeven}{-0.18}
\newcommand{\EThreeGPFreqHalfCanaryAware}{3.77}
\newcommand{\EThreeGPFreqHalfCanaryGate}{7.67}
\newcommand{\EThreeGPFreqHalfCanaryRG}{3.06}
\newcommand{\EThreeGPFreqHalfChecks}{6.0}
\newcommand{\EThreeGPFreqHalfGFifth}{11.59}
\newcommand{\EThreeGPFreqHalfGFifthD}{2.31}
\newcommand{\EThreeGPFreqHalfGFifthHi}{3.48}
\newcommand{\EThreeGPFreqHalfGFifthLo}{1.28}
\newcommand{\EThreeGPFreqHalfGHalf}{13.40}
\newcommand{\EThreeGPFreqHalfGHalfD}{4.12}
\newcommand{\EThreeGPFreqHalfGHalfHi}{5.32}
\newcommand{\EThreeGPFreqHalfGHalfLo}{3.06}
\newcommand{\EThreeGPFreqHalfGOne}{16.43}
\newcommand{\EThreeGPFreqHalfGOneD}{7.15}
\newcommand{\EThreeGPFreqHalfGOneHi}{8.42}
\newcommand{\EThreeGPFreqHalfGOneLo}{6.01}
\newcommand{\EThreeGPFreqHalfGTenth}{10.98}
\newcommand{\EThreeGPFreqHalfGTenthD}{1.70}
\newcommand{\EThreeGPFreqHalfGTenthHi}{2.86}
\newcommand{\EThreeGPFreqHalfGTenthLo}{0.68}
\newcommand{\EThreeGPFreqHalfUFifth}{10.16}
\newcommand{\EThreeGPFreqHalfUFifthD}{0.88}
\newcommand{\EThreeGPFreqHalfUFifthHi}{1.35}
\newcommand{\EThreeGPFreqHalfUFifthLo}{0.38}
\newcommand{\EThreeGPFreqHalfUHalf}{11.95}
\newcommand{\EThreeGPFreqHalfUHalfD}{2.67}
\newcommand{\EThreeGPFreqHalfUHalfHi}{3.16}
\newcommand{\EThreeGPFreqHalfUHalfLo}{2.15}
\newcommand{\EThreeGPFreqHalfUOne}{14.94}
\newcommand{\EThreeGPFreqHalfUOneD}{5.66}
\newcommand{\EThreeGPFreqHalfUOneHi}{6.23}
\newcommand{\EThreeGPFreqHalfUOneLo}{5.08}
\newcommand{\EThreeGPFreqHalfUTenth}{9.56}
\newcommand{\EThreeGPFreqHalfUTenthD}{0.28}
\newcommand{\EThreeGPFreqHalfUTenthHi}{0.75}
\newcommand{\EThreeGPFreqHalfUTenthLo}{-0.22}
\newcommand{\EThreeGPRareFullAFifth}{7.77}
\newcommand{\EThreeGPRareFullAFifthD}{-1.51}
\newcommand{\EThreeGPRareFullAFifthHi}{-0.67}
\newcommand{\EThreeGPRareFullAFifthLo}{-2.42}
\newcommand{\EThreeGPRareFullAHalf}{8.10}
\newcommand{\EThreeGPRareFullAHalfD}{-1.18}
\newcommand{\EThreeGPRareFullAHalfHi}{-0.33}
\newcommand{\EThreeGPRareFullAHalfLo}{-2.08}
\newcommand{\EThreeGPRareFullAOne}{8.66}
\newcommand{\EThreeGPRareFullAOneD}{-0.62}
\newcommand{\EThreeGPRareFullAOneHi}{0.25}
\newcommand{\EThreeGPRareFullAOneLo}{-1.51}
\newcommand{\EThreeGPRareFullATenth}{7.65}
\newcommand{\EThreeGPRareFullATenthD}{-1.63}
\newcommand{\EThreeGPRareFullATenthHi}{-0.79}
\newcommand{\EThreeGPRareFullATenthLo}{-2.53}
\newcommand{\EThreeGPRareFullAmGFifth}{-0.21}
\newcommand{\EThreeGPRareFullAmGFifthHi}{0.10}
\newcommand{\EThreeGPRareFullAmGFifthLo}{-0.56}
\newcommand{\EThreeGPRareFullAmGHalf}{-0.22}
\newcommand{\EThreeGPRareFullAmGHalfHi}{0.09}
\newcommand{\EThreeGPRareFullAmGHalfLo}{-0.57}
\newcommand{\EThreeGPRareFullAmGOne}{-0.24}
\newcommand{\EThreeGPRareFullAmGOneHi}{0.07}
\newcommand{\EThreeGPRareFullAmGOneLo}{-0.60}
\newcommand{\EThreeGPRareFullAmGTenth}{-0.21}
\newcommand{\EThreeGPRareFullAmGTenthHi}{0.10}
\newcommand{\EThreeGPRareFullAmGTenthLo}{-0.55}
\newcommand{\EThreeGPRareFullBreakeven}{1.32}
\newcommand{\EThreeGPRareFullCanaryAware}{0.71}
\newcommand{\EThreeGPRareFullCanaryGate}{1.49}
\newcommand{\EThreeGPRareFullCanaryRG}{0.42}
\newcommand{\EThreeGPRareFullChecks}{1.2}
\newcommand{\EThreeGPRareFullGFifth}{7.98}
\newcommand{\EThreeGPRareFullGFifthD}{-1.30}
\newcommand{\EThreeGPRareFullGFifthHi}{-0.29}
\newcommand{\EThreeGPRareFullGFifthLo}{-2.33}
\newcommand{\EThreeGPRareFullGHalf}{8.32}
\newcommand{\EThreeGPRareFullGHalfD}{-0.95}
\newcommand{\EThreeGPRareFullGHalfHi}{0.07}
\newcommand{\EThreeGPRareFullGHalfLo}{-2.01}
\newcommand{\EThreeGPRareFullGOne}{8.90}
\newcommand{\EThreeGPRareFullGOneD}{-0.38}
\newcommand{\EThreeGPRareFullGOneHi}{0.67}
\newcommand{\EThreeGPRareFullGOneLo}{-1.44}
\newcommand{\EThreeGPRareFullGTenth}{7.86}
\newcommand{\EThreeGPRareFullGTenthD}{-1.42}
\newcommand{\EThreeGPRareFullGTenthHi}{-0.41}
\newcommand{\EThreeGPRareFullGTenthLo}{-2.44}
\newcommand{\EThreeGPRareFullUFifth}{7.70}
\newcommand{\EThreeGPRareFullUFifthD}{-1.58}
\newcommand{\EThreeGPRareFullUFifthHi}{-0.87}
\newcommand{\EThreeGPRareFullUFifthLo}{-2.39}
\newcommand{\EThreeGPRareFullUHalf}{8.01}
\newcommand{\EThreeGPRareFullUHalfD}{-1.27}
\newcommand{\EThreeGPRareFullUHalfHi}{-0.56}
\newcommand{\EThreeGPRareFullUHalfLo}{-2.06}
\newcommand{\EThreeGPRareFullUOne}{8.52}
\newcommand{\EThreeGPRareFullUOneD}{-0.76}
\newcommand{\EThreeGPRareFullUOneHi}{-0.07}
\newcommand{\EThreeGPRareFullUOneLo}{-1.53}
\newcommand{\EThreeGPRareFullUTenth}{7.60}
\newcommand{\EThreeGPRareFullUTenthD}{-1.68}
\newcommand{\EThreeGPRareFullUTenthHi}{-0.97}
\newcommand{\EThreeGPRareFullUTenthLo}{-2.50}
\newcommand{\EThreeGPRareHalfAFifth}{8.83}
\newcommand{\EThreeGPRareHalfAFifthD}{-0.45}
\newcommand{\EThreeGPRareHalfAFifthHi}{0.18}
\newcommand{\EThreeGPRareHalfAFifthLo}{-1.09}
\newcommand{\EThreeGPRareHalfAHalf}{9.12}
\newcommand{\EThreeGPRareHalfAHalfD}{-0.15}
\newcommand{\EThreeGPRareHalfAHalfHi}{0.46}
\newcommand{\EThreeGPRareHalfAHalfLo}{-0.78}
\newcommand{\EThreeGPRareHalfAOne}{9.62}
\newcommand{\EThreeGPRareHalfAOneD}{0.34}
\newcommand{\EThreeGPRareHalfAOneHi}{0.94}
\newcommand{\EThreeGPRareHalfAOneLo}{-0.28}
\newcommand{\EThreeGPRareHalfATenth}{8.73}
\newcommand{\EThreeGPRareHalfATenthD}{-0.55}
\newcommand{\EThreeGPRareHalfATenthHi}{0.09}
\newcommand{\EThreeGPRareHalfATenthLo}{-1.19}
\newcommand{\EThreeGPRareHalfAmGFifth}{-1.29}
\newcommand{\EThreeGPRareHalfAmGFifthHi}{-0.53}
\newcommand{\EThreeGPRareHalfAmGFifthLo}{-2.15}
\newcommand{\EThreeGPRareHalfAmGHalf}{-1.30}
\newcommand{\EThreeGPRareHalfAmGHalfHi}{-0.54}
\newcommand{\EThreeGPRareHalfAmGHalfLo}{-2.17}
\newcommand{\EThreeGPRareHalfAmGOne}{-1.33}
\newcommand{\EThreeGPRareHalfAmGOneHi}{-0.57}
\newcommand{\EThreeGPRareHalfAmGOneLo}{-2.20}
\newcommand{\EThreeGPRareHalfAmGTenth}{-1.28}
\newcommand{\EThreeGPRareHalfAmGTenthHi}{-0.52}
\newcommand{\EThreeGPRareHalfAmGTenthLo}{-2.14}
\newcommand{\EThreeGPRareHalfBreakeven}{-0.61}
\newcommand{\EThreeGPRareHalfCanaryAware}{0.69}
\newcommand{\EThreeGPRareHalfCanaryGate}{1.76}
\newcommand{\EThreeGPRareHalfCanaryRG}{0.55}
\newcommand{\EThreeGPRareHalfChecks}{1.0}
\newcommand{\EThreeGPRareHalfGFifth}{10.12}
\newcommand{\EThreeGPRareHalfGFifthD}{0.84}
\newcommand{\EThreeGPRareHalfGFifthHi}{2.01}
\newcommand{\EThreeGPRareHalfGFifthLo}{-0.16}
\newcommand{\EThreeGPRareHalfGHalf}{10.43}
\newcommand{\EThreeGPRareHalfGHalfD}{1.15}
\newcommand{\EThreeGPRareHalfGHalfHi}{2.32}
\newcommand{\EThreeGPRareHalfGHalfLo}{0.15}
\newcommand{\EThreeGPRareHalfGOne}{10.95}
\newcommand{\EThreeGPRareHalfGOneD}{1.67}
\newcommand{\EThreeGPRareHalfGOneHi}{2.84}
\newcommand{\EThreeGPRareHalfGOneLo}{0.64}
\newcommand{\EThreeGPRareHalfGTenth}{10.02}
\newcommand{\EThreeGPRareHalfGTenthD}{0.74}
\newcommand{\EThreeGPRareHalfGTenthHi}{1.91}
\newcommand{\EThreeGPRareHalfGTenthLo}{-0.26}
\newcommand{\EThreeGPRareHalfUFifth}{8.69}
\newcommand{\EThreeGPRareHalfUFifthD}{-0.59}
\newcommand{\EThreeGPRareHalfUFifthHi}{-0.13}
\newcommand{\EThreeGPRareHalfUFifthLo}{-1.09}
\newcommand{\EThreeGPRareHalfUHalf}{8.98}
\newcommand{\EThreeGPRareHalfUHalfD}{-0.30}
\newcommand{\EThreeGPRareHalfUHalfHi}{0.15}
\newcommand{\EThreeGPRareHalfUHalfLo}{-0.79}
\newcommand{\EThreeGPRareHalfUOne}{9.46}
\newcommand{\EThreeGPRareHalfUOneD}{0.18}
\newcommand{\EThreeGPRareHalfUOneHi}{0.62}
\newcommand{\EThreeGPRareHalfUOneLo}{-0.30}
\newcommand{\EThreeGPRareHalfUTenth}{8.59}
\newcommand{\EThreeGPRareHalfUTenthD}{-0.69}
\newcommand{\EThreeGPRareHalfUTenthHi}{-0.22}
\newcommand{\EThreeGPRareHalfUTenthLo}{-1.18}
\newcommand{\EThreeGateFreqFullD}{-0.38}
\newcommand{\EThreeGateFreqFullDhi}{0.27}
\newcommand{\EThreeGateFreqFullDlo}{-1.18}
\newcommand{\EThreeGateFreqFullGR}{5.10}
\newcommand{\EThreeGateFreqFullGRhi}{6.21}
\newcommand{\EThreeGateFreqFullGRlo}{4.10}
\newcommand{\EThreeGateFreqFullGate}{14.38}
\newcommand{\EThreeGateFreqFullRG}{14.00}
\newcommand{\EThreeGateFreqHalfD}{-1.48}
\newcommand{\EThreeGateFreqHalfDhi}{-0.65}
\newcommand{\EThreeGateFreqHalfDlo}{-2.50}
\newcommand{\EThreeGateFreqHalfGR}{7.15}
\newcommand{\EThreeGateFreqHalfGRhi}{8.49}
\newcommand{\EThreeGateFreqHalfGRlo}{5.92}
\newcommand{\EThreeGateFreqHalfGate}{16.43}
\newcommand{\EThreeGateFreqHalfRG}{14.94}
\newcommand{\EThreeGateRareFullD}{-0.38}
\newcommand{\EThreeGateRareFullDhi}{0.27}
\newcommand{\EThreeGateRareFullDlo}{-1.18}
\newcommand{\EThreeGateRareFullGR}{-0.38}
\newcommand{\EThreeGateRareFullGRhi}{0.69}
\newcommand{\EThreeGateRareFullGRlo}{-1.35}
\newcommand{\EThreeGateRareFullGate}{8.90}
\newcommand{\EThreeGateRareFullRG}{8.52}
\newcommand{\EThreeGateRareHalfD}{-1.48}
\newcommand{\EThreeGateRareHalfDhi}{-0.65}
\newcommand{\EThreeGateRareHalfDlo}{-2.50}
\newcommand{\EThreeGateRareHalfGR}{1.67}
\newcommand{\EThreeGateRareHalfGRhi}{2.88}
\newcommand{\EThreeGateRareHalfGRlo}{0.58}
\newcommand{\EThreeGateRareHalfGate}{10.95}
\newcommand{\EThreeGateRareHalfRG}{9.46}
\newcommand{\EThreeGitCalls}{1{,}669{,}886}
\newcommand{\EThreeHTwoIrrOne}{5.4}
\newcommand{\EThreeHTwoIrrOneHi}{8.3}
\newcommand{\EThreeHTwoIrrOneLo}{2.9}
\newcommand{\EThreeHTwoIrrOneP}{=0.002}
\newcommand{\EThreeHTwoIrrTen}{5.0}
\newcommand{\EThreeHTwoIrrTenHi}{7.5}
\newcommand{\EThreeHTwoIrrTenLo}{2.5}
\newcommand{\EThreeHTwoIrrTenP}{=0.001}
\newcommand{\EThreeHTwoIrrTenth}{4.2}
\newcommand{\EThreeHTwoIrrTenthHi}{6.7}
\newcommand{\EThreeHTwoIrrTenthLo}{1.7}
\newcommand{\EThreeHTwoIrrTenthP}{=0.004}
\newcommand{\EThreeHTwoIrrZero}{3.8}
\newcommand{\EThreeHTwoIrrZeroHi}{6.2}
\newcommand{\EThreeHTwoIrrZeroLo}{1.2}
\newcommand{\EThreeHTwoIrrZeroP}{=0.007}
\newcommand{\EThreeIrrPerBl}{21.2}
\newcommand{\EThreeIrrPerCo}{11.2}
\newcommand{\EThreeIrrPerMy}{20.4}
\newcommand{\EThreeIrrPerNa}{15.8}
\newcommand{\EThreeIrrPerNo}{23.8}
\newcommand{\EThreeIrrPerPe}{8.8}
\newcommand{\EThreeIrrPerRb}{15.0}
\newcommand{\EThreeIrrPerRover}{10.4}
\newcommand{\EThreeIrrPerVa}{6.7}
\newcommand{\EThreeJWCo}{11.9}
\newcommand{\EThreeJWCoC}{skill, late}
\newcommand{\EThreeJWMy}{12.7}
\newcommand{\EThreeJWMyC}{skill}
\newcommand{\EThreeJWNa}{12.7}
\newcommand{\EThreeJWNaC}{skill, late}
\newcommand{\EThreeJWNo}{14.3}
\newcommand{\EThreeJWNoC}{skill}
\newcommand{\EThreeJWRb}{14.8}
\newcommand{\EThreeJWRbC}{memory note, late}
\newcommand{\EThreeJWRover}{12.6}
\newcommand{\EThreeJWRoverC}{skill, late}
\newcommand{\EThreeKDAlAtt}{64}
\newcommand{\EThreeKDAlMy}{1.10}
\newcommand{\EThreeKDAlMyHi}{2.13}
\newcommand{\EThreeKDAlMyIrr}{6.7}
\newcommand{\EThreeKDAlMyIrrHi}{10.0}
\newcommand{\EThreeKDAlMyIrrLo}{3.3}
\newcommand{\EThreeKDAlMyIrrP}{=0.001}
\newcommand{\EThreeKDAlMyLo}{0.17}
\newcommand{\EThreeKDAlMyLoses}{13}
\newcommand{\EThreeKDAlMyP}{=0.070}
\newcommand{\EThreeKDAlMyWins}{17}
\newcommand{\EThreeKDAlNo}{1.39}
\newcommand{\EThreeKDAlNoHi}{2.35}
\newcommand{\EThreeKDAlNoIrr}{9.6}
\newcommand{\EThreeKDAlNoIrrHi}{12.9}
\newcommand{\EThreeKDAlNoIrrLo}{6.2}
\newcommand{\EThreeKDAlNoIrrP}{<0.001}
\newcommand{\EThreeKDAlNoLo}{0.46}
\newcommand{\EThreeKDAlNoLoses}{10}
\newcommand{\EThreeKDAlNoP}{=0.022}
\newcommand{\EThreeKDAlNoWins}{20}
\newcommand{\EThreeKDAlRb}{0.89}
\newcommand{\EThreeKDAlRbHi}{1.73}
\newcommand{\EThreeKDAlRbIrr}{0.8}
\newcommand{\EThreeKDAlRbIrrHi}{3.8}
\newcommand{\EThreeKDAlRbIrrLo}{-1.7}
\newcommand{\EThreeKDAlRbIrrP}{=0.564}
\newcommand{\EThreeKDAlRbLo}{0.10}
\newcommand{\EThreeKDAlRbLoses}{11}
\newcommand{\EThreeKDAlRbP}{=0.082}
\newcommand{\EThreeKDAlRbWins}{19}
\newcommand{\EThreeKDAlRover}{10.47}
\newcommand{\EThreeKDAlRoverIrr}{14.2}
\newcommand{\EThreeKRFiveCalAtt}{5}
\newcommand{\EThreeKRFiveCalNa}{-0.04}
\newcommand{\EThreeKRFiveCalNaHi}{-0.01}
\newcommand{\EThreeKRFiveCalNaIrr}{0.0}
\newcommand{\EThreeKRFiveCalNaIrrHi}{0.0}
\newcommand{\EThreeKRFiveCalNaIrrLo}{0.0}
\newcommand{\EThreeKRFiveCalNaIrrP}{=1.000}
\newcommand{\EThreeKRFiveCalNaLo}{-0.07}
\newcommand{\EThreeKRFiveCalNaLoses}{8}
\newcommand{\EThreeKRFiveCalNaP}{=0.021}
\newcommand{\EThreeKRFiveCalNaWins}{1}
\newcommand{\EThreeKRFiveCalNo}{0.25}
\newcommand{\EThreeKRFiveCalNoHi}{0.59}
\newcommand{\EThreeKRFiveCalNoIrr}{1.2}
\newcommand{\EThreeKRFiveCalNoIrrHi}{2.9}
\newcommand{\EThreeKRFiveCalNoIrrLo}{0.0}
\newcommand{\EThreeKRFiveCalNoIrrP}{=0.083}
\newcommand{\EThreeKRFiveCalNoLo}{-0.03}
\newcommand{\EThreeKRFiveCalNoLoses}{6}
\newcommand{\EThreeKRFiveCalNoP}{=0.854}
\newcommand{\EThreeKRFiveCalNoWins}{3}
\newcommand{\EThreeKRFiveCalRb}{0.10}
\newcommand{\EThreeKRFiveCalRbHi}{0.30}
\newcommand{\EThreeKRFiveCalRbIrr}{0.4}
\newcommand{\EThreeKRFiveCalRbIrrHi}{1.2}
\newcommand{\EThreeKRFiveCalRbIrrLo}{0.0}
\newcommand{\EThreeKRFiveCalRbIrrP}{=0.317}
\newcommand{\EThreeKRFiveCalRbLo}{-0.04}
\newcommand{\EThreeKRFiveCalRbLoses}{6}
\newcommand{\EThreeKRFiveCalRbP}{=0.807}
\newcommand{\EThreeKRFiveCalRbWins}{3}
\newcommand{\EThreeKRFiveCalRover}{0.92}
\newcommand{\EThreeKRFiveCalRoverIrr}{1.7}
\newcommand{\EThreeKRFiveFixAtt}{5}
\newcommand{\EThreeKRFiveFixNa}{-0.17}
\newcommand{\EThreeKRFiveFixNaHi}{0.20}
\newcommand{\EThreeKRFiveFixNaIrr}{0.8}
\newcommand{\EThreeKRFiveFixNaIrrHi}{2.1}
\newcommand{\EThreeKRFiveFixNaIrrLo}{0.0}
\newcommand{\EThreeKRFiveFixNaIrrP}{=0.157}
\newcommand{\EThreeKRFiveFixNaLo}{-0.57}
\newcommand{\EThreeKRFiveFixNaLoses}{19}
\newcommand{\EThreeKRFiveFixNaP}{=0.381}
\newcommand{\EThreeKRFiveFixNaWins}{10}
\newcommand{\EThreeKRFiveFixNo}{-0.16}
\newcommand{\EThreeKRFiveFixNoHi}{0.30}
\newcommand{\EThreeKRFiveFixNoIrr}{2.9}
\newcommand{\EThreeKRFiveFixNoIrrHi}{5.0}
\newcommand{\EThreeKRFiveFixNoIrrLo}{1.2}
\newcommand{\EThreeKRFiveFixNoIrrP}{=0.008}
\newcommand{\EThreeKRFiveFixNoLo}{-0.60}
\newcommand{\EThreeKRFiveFixNoLoses}{21}
\newcommand{\EThreeKRFiveFixNoP}{=0.224}
\newcommand{\EThreeKRFiveFixNoWins}{9}
\newcommand{\EThreeKRFiveFixRb}{0.68}
\newcommand{\EThreeKRFiveFixRbHi}{1.30}
\newcommand{\EThreeKRFiveFixRbIrr}{1.2}
\newcommand{\EThreeKRFiveFixRbIrrHi}{2.9}
\newcommand{\EThreeKRFiveFixRbIrrLo}{0.0}
\newcommand{\EThreeKRFiveFixRbIrrP}{=0.083}
\newcommand{\EThreeKRFiveFixRbLo}{0.09}
\newcommand{\EThreeKRFiveFixRbLoses}{12}
\newcommand{\EThreeKRFiveFixRbP}{=0.029}
\newcommand{\EThreeKRFiveFixRbWins}{17}
\newcommand{\EThreeKRFiveFixRover}{1.32}
\newcommand{\EThreeKRFiveFixRoverIrr}{0.0}
\newcommand{\EThreeKRTwentyCalAtt}{18}
\newcommand{\EThreeKRTwentyCalNa}{0.23}
\newcommand{\EThreeKRTwentyCalNaHi}{0.55}
\newcommand{\EThreeKRTwentyCalNaIrr}{0.8}
\newcommand{\EThreeKRTwentyCalNaIrrHi}{2.1}
\newcommand{\EThreeKRTwentyCalNaIrrLo}{0.0}
\newcommand{\EThreeKRTwentyCalNaIrrP}{=0.157}
\newcommand{\EThreeKRTwentyCalNaLo}{-0.07}
\newcommand{\EThreeKRTwentyCalNaLoses}{16}
\newcommand{\EThreeKRTwentyCalNaP}{=0.888}
\newcommand{\EThreeKRTwentyCalNaWins}{10}
\newcommand{\EThreeKRTwentyCalNo}{0.60}
\newcommand{\EThreeKRTwentyCalNoHi}{1.11}
\newcommand{\EThreeKRTwentyCalNoIrr}{3.3}
\newcommand{\EThreeKRTwentyCalNoIrrHi}{5.4}
\newcommand{\EThreeKRTwentyCalNoIrrLo}{1.2}
\newcommand{\EThreeKRTwentyCalNoIrrP}{=0.005}
\newcommand{\EThreeKRTwentyCalNoLo}{0.10}
\newcommand{\EThreeKRTwentyCalNoLoses}{16}
\newcommand{\EThreeKRTwentyCalNoP}{=0.444}
\newcommand{\EThreeKRTwentyCalNoWins}{10}
\newcommand{\EThreeKRTwentyCalRb}{0.40}
\newcommand{\EThreeKRTwentyCalRbHi}{0.75}
\newcommand{\EThreeKRTwentyCalRbIrr}{0.8}
\newcommand{\EThreeKRTwentyCalRbIrrHi}{2.1}
\newcommand{\EThreeKRTwentyCalRbIrrLo}{0.0}
\newcommand{\EThreeKRTwentyCalRbIrrP}{=0.157}
\newcommand{\EThreeKRTwentyCalRbLo}{0.10}
\newcommand{\EThreeKRTwentyCalRbLoses}{14}
\newcommand{\EThreeKRTwentyCalRbP}{=0.301}
\newcommand{\EThreeKRTwentyCalRbWins}{12}
\newcommand{\EThreeKRTwentyCalRover}{2.65}
\newcommand{\EThreeKRTwentyCalRoverIrr}{2.9}
\newcommand{\EThreeKRTwentyFixAtt}{18}
\newcommand{\EThreeKRTwentyFixNa}{-0.07}
\newcommand{\EThreeKRTwentyFixNaHi}{0.35}
\newcommand{\EThreeKRTwentyFixNaIrr}{0.8}
\newcommand{\EThreeKRTwentyFixNaIrrHi}{2.1}
\newcommand{\EThreeKRTwentyFixNaIrrLo}{0.0}
\newcommand{\EThreeKRTwentyFixNaIrrP}{=0.157}
\newcommand{\EThreeKRTwentyFixNaLo}{-0.50}
\newcommand{\EThreeKRTwentyFixNaLoses}{16}
\newcommand{\EThreeKRTwentyFixNaP}{=0.897}
\newcommand{\EThreeKRTwentyFixNaWins}{13}
\newcommand{\EThreeKRTwentyFixNo}{0.20}
\newcommand{\EThreeKRTwentyFixNoHi}{0.78}
\newcommand{\EThreeKRTwentyFixNoIrr}{4.2}
\newcommand{\EThreeKRTwentyFixNoIrrHi}{6.2}
\newcommand{\EThreeKRTwentyFixNoIrrLo}{2.1}
\newcommand{\EThreeKRTwentyFixNoIrrP}{=0.002}
\newcommand{\EThreeKRTwentyFixNoLo}{-0.38}
\newcommand{\EThreeKRTwentyFixNoLoses}{17}
\newcommand{\EThreeKRTwentyFixNoP}{=0.934}
\newcommand{\EThreeKRTwentyFixNoWins}{13}
\newcommand{\EThreeKRTwentyFixRb}{0.80}
\newcommand{\EThreeKRTwentyFixRbHi}{1.45}
\newcommand{\EThreeKRTwentyFixRbIrr}{1.2}
\newcommand{\EThreeKRTwentyFixRbIrrHi}{2.9}
\newcommand{\EThreeKRTwentyFixRbIrrLo}{0.0}
\newcommand{\EThreeKRTwentyFixRbIrrP}{=0.083}
\newcommand{\EThreeKRTwentyFixRbLo}{0.21}
\newcommand{\EThreeKRTwentyFixRbLoses}{10}
\newcommand{\EThreeKRTwentyFixRbP}{=0.011}
\newcommand{\EThreeKRTwentyFixRbWins}{20}
\newcommand{\EThreeKRTwentyFixRover}{3.05}
\newcommand{\EThreeKRTwentyFixRoverIrr}{2.1}
\newcommand{\EThreeKSEightAtt}{64}
\newcommand{\EThreeKSEightCo}{-0.32}
\newcommand{\EThreeKSEightCoHi}{0.20}
\newcommand{\EThreeKSEightCoIrr}{-2.5}
\newcommand{\EThreeKSEightCoIrrHi}{-0.8}
\newcommand{\EThreeKSEightCoIrrLo}{-5.0}
\newcommand{\EThreeKSEightCoIrrP}{=0.034}
\newcommand{\EThreeKSEightCoLo}{-0.88}
\newcommand{\EThreeKSEightCoLoses}{12}
\newcommand{\EThreeKSEightCoP}{=0.505}
\newcommand{\EThreeKSEightCoWins}{10}
\newcommand{\EThreeKSEightMy}{-0.46}
\newcommand{\EThreeKSEightMyHi}{0.29}
\newcommand{\EThreeKSEightMyIrr}{2.1}
\newcommand{\EThreeKSEightMyIrrHi}{4.6}
\newcommand{\EThreeKSEightMyIrrLo}{-0.0}
\newcommand{\EThreeKSEightMyIrrP}{=0.096}
\newcommand{\EThreeKSEightMyLo}{-1.14}
\newcommand{\EThreeKSEightMyLoses}{22}
\newcommand{\EThreeKSEightMyP}{=0.150}
\newcommand{\EThreeKSEightMyWins}{8}
\newcommand{\EThreeKSEightNa}{-0.40}
\newcommand{\EThreeKSEightNaHi}{0.03}
\newcommand{\EThreeKSEightNaIrr}{1.2}
\newcommand{\EThreeKSEightNaIrrHi}{2.5}
\newcommand{\EThreeKSEightNaIrrLo}{0.0}
\newcommand{\EThreeKSEightNaIrrP}{=0.083}
\newcommand{\EThreeKSEightNaLo}{-0.80}
\newcommand{\EThreeKSEightNaLoses}{21}
\newcommand{\EThreeKSEightNaP}{=0.030}
\newcommand{\EThreeKSEightNaWins}{6}
\newcommand{\EThreeKSEightNo}{-0.53}
\newcommand{\EThreeKSEightNoHi}{0.11}
\newcommand{\EThreeKSEightNoIrr}{2.9}
\newcommand{\EThreeKSEightNoIrrHi}{5.0}
\newcommand{\EThreeKSEightNoIrrLo}{1.2}
\newcommand{\EThreeKSEightNoIrrP}{=0.008}
\newcommand{\EThreeKSEightNoLo}{-1.10}
\newcommand{\EThreeKSEightNoLoses}{23}
\newcommand{\EThreeKSEightNoP}{=0.100}
\newcommand{\EThreeKSEightNoWins}{7}
\newcommand{\EThreeKSEightRb}{1.45}
\newcommand{\EThreeKSEightRbHi}{2.17}
\newcommand{\EThreeKSEightRbIrr}{-0.4}
\newcommand{\EThreeKSEightRbIrrHi}{1.2}
\newcommand{\EThreeKSEightRbIrrLo}{-2.5}
\newcommand{\EThreeKSEightRbIrrP}{=0.705}
\newcommand{\EThreeKSEightRbLo}{0.68}
\newcommand{\EThreeKSEightRbLoses}{6}
\newcommand{\EThreeKSEightRbP}{=0.002}
\newcommand{\EThreeKSEightRbWins}{24}
\newcommand{\EThreeKSEightRover}{10.59}
\newcommand{\EThreeKSEightRoverIrr}{13.8}
\newcommand{\EThreeKSFourAtt}{64}
\newcommand{\EThreeKSFourCo}{0.16}
\newcommand{\EThreeKSFourCoHi}{0.62}
\newcommand{\EThreeKSFourCoIrr}{-0.8}
\newcommand{\EThreeKSFourCoIrrHi}{1.2}
\newcommand{\EThreeKSFourCoIrrLo}{-2.9}
\newcommand{\EThreeKSFourCoIrrP}{=0.479}
\newcommand{\EThreeKSFourCoLo}{-0.37}
\newcommand{\EThreeKSFourCoLoses}{10}
\newcommand{\EThreeKSFourCoP}{=0.182}
\newcommand{\EThreeKSFourCoWins}{14}
\newcommand{\EThreeKSFourNa}{0.33}
\newcommand{\EThreeKSFourNaHi}{0.88}
\newcommand{\EThreeKSFourNaIrr}{3.8}
\newcommand{\EThreeKSFourNaIrrHi}{6.2}
\newcommand{\EThreeKSFourNaIrrLo}{1.2}
\newcommand{\EThreeKSFourNaIrrP}{=0.014}
\newcommand{\EThreeKSFourNaLo}{-0.25}
\newcommand{\EThreeKSFourNaLoses}{16}
\newcommand{\EThreeKSFourNaP}{=0.387}
\newcommand{\EThreeKSFourNaWins}{14}
\newcommand{\EThreeKSFourNo}{0.80}
\newcommand{\EThreeKSFourNoHi}{1.66}
\newcommand{\EThreeKSFourNoIrr}{7.9}
\newcommand{\EThreeKSFourNoIrrHi}{11.2}
\newcommand{\EThreeKSFourNoIrrLo}{4.6}
\newcommand{\EThreeKSFourNoIrrP}{<0.001}
\newcommand{\EThreeKSFourNoLo}{-0.03}
\newcommand{\EThreeKSFourNoLoses}{13}
\newcommand{\EThreeKSFourNoP}{=0.101}
\newcommand{\EThreeKSFourNoWins}{15}
\newcommand{\EThreeKSFourRb}{1.50}
\newcommand{\EThreeKSFourRbHi}{2.41}
\newcommand{\EThreeKSFourRbIrr}{0.8}
\newcommand{\EThreeKSFourRbIrrHi}{3.8}
\newcommand{\EThreeKSFourRbIrrLo}{-2.5}
\newcommand{\EThreeKSFourRbIrrP}{=0.627}
\newcommand{\EThreeKSFourRbLo}{0.57}
\newcommand{\EThreeKSFourRbLoses}{7}
\newcommand{\EThreeKSFourRbP}{=0.003}
\newcommand{\EThreeKSFourRbWins}{22}
\newcommand{\EThreeKSFourRover}{10.60}
\newcommand{\EThreeKSFourRoverIrr}{12.9}
\newcommand{\EThreeKSOneLowAtt}{64}
\newcommand{\EThreeKSOneLowMy}{0.86}
\newcommand{\EThreeKSOneLowMyHi}{1.83}
\newcommand{\EThreeKSOneLowMyIrr}{5.8}
\newcommand{\EThreeKSOneLowMyIrrHi}{9.6}
\newcommand{\EThreeKSOneLowMyIrrLo}{2.9}
\newcommand{\EThreeKSOneLowMyIrrP}{=0.004}
\newcommand{\EThreeKSOneLowMyLo}{-0.02}
\newcommand{\EThreeKSOneLowMyLoses}{17}
\newcommand{\EThreeKSOneLowMyP}{=0.376}
\newcommand{\EThreeKSOneLowMyWins}{13}
\newcommand{\EThreeKSOneLowNo}{0.96}
\newcommand{\EThreeKSOneLowNoHi}{2.00}
\newcommand{\EThreeKSOneLowNoIrr}{6.7}
\newcommand{\EThreeKSOneLowNoIrrHi}{10.4}
\newcommand{\EThreeKSOneLowNoIrrLo}{3.3}
\newcommand{\EThreeKSOneLowNoIrrP}{=0.002}
\newcommand{\EThreeKSOneLowNoLo}{0.03}
\newcommand{\EThreeKSOneLowNoLoses}{17}
\newcommand{\EThreeKSOneLowNoP}{=0.349}
\newcommand{\EThreeKSOneLowNoWins}{13}
\newcommand{\EThreeKSOneLowRb}{1.59}
\newcommand{\EThreeKSOneLowRbHi}{2.61}
\newcommand{\EThreeKSOneLowRbIrr}{3.8}
\newcommand{\EThreeKSOneLowRbIrrHi}{6.7}
\newcommand{\EThreeKSOneLowRbIrrLo}{1.2}
\newcommand{\EThreeKSOneLowRbIrrP}{=0.024}
\newcommand{\EThreeKSOneLowRbLo}{0.72}
\newcommand{\EThreeKSOneLowRbLoses}{9}
\newcommand{\EThreeKSOneLowRbP}{=0.001}
\newcommand{\EThreeKSOneLowRbWins}{21}
\newcommand{\EThreeKSOneLowRover}{10.90}
\newcommand{\EThreeKSOneLowRoverIrr}{17.1}
\newcommand{\EThreeKSSixAtt}{64}
\newcommand{\EThreeKSSixCo}{-0.16}
\newcommand{\EThreeKSSixCoHi}{0.33}
\newcommand{\EThreeKSSixCoIrr}{-2.1}
\newcommand{\EThreeKSSixCoIrrHi}{0.0}
\newcommand{\EThreeKSSixCoIrrLo}{-4.6}
\newcommand{\EThreeKSSixCoIrrP}{=0.096}
\newcommand{\EThreeKSSixCoLo}{-0.70}
\newcommand{\EThreeKSSixCoLoses}{12}
\newcommand{\EThreeKSSixCoP}{=0.782}
\newcommand{\EThreeKSSixCoWins}{10}
\newcommand{\EThreeKSSixNa}{-0.01}
\newcommand{\EThreeKSSixNaHi}{0.52}
\newcommand{\EThreeKSSixNaIrr}{2.5}
\newcommand{\EThreeKSSixNaIrrHi}{5.0}
\newcommand{\EThreeKSSixNaIrrLo}{0.4}
\newcommand{\EThreeKSSixNaIrrP}{=0.034}
\newcommand{\EThreeKSSixNaLo}{-0.49}
\newcommand{\EThreeKSSixNaLoses}{18}
\newcommand{\EThreeKSSixNaP}{=0.577}
\newcommand{\EThreeKSSixNaWins}{10}
\newcommand{\EThreeKSSixNo}{-0.09}
\newcommand{\EThreeKSSixNoHi}{0.60}
\newcommand{\EThreeKSSixNoIrr}{5.0}
\newcommand{\EThreeKSSixNoIrrHi}{7.9}
\newcommand{\EThreeKSSixNoIrrLo}{2.5}
\newcommand{\EThreeKSSixNoIrrP}{=0.001}
\newcommand{\EThreeKSSixNoLo}{-0.72}
\newcommand{\EThreeKSSixNoLoses}{19}
\newcommand{\EThreeKSSixNoP}{=0.713}
\newcommand{\EThreeKSSixNoWins}{10}
\newcommand{\EThreeKSSixRb}{1.41}
\newcommand{\EThreeKSSixRbHi}{2.27}
\newcommand{\EThreeKSSixRbIrr}{-0.8}
\newcommand{\EThreeKSSixRbIrrHi}{1.2}
\newcommand{\EThreeKSSixRbIrrLo}{-3.3}
\newcommand{\EThreeKSSixRbIrrP}{=0.479}
\newcommand{\EThreeKSSixRbLo}{0.56}
\newcommand{\EThreeKSSixRbLoses}{7}
\newcommand{\EThreeKSSixRbP}{=0.005}
\newcommand{\EThreeKSSixRbWins}{23}
\newcommand{\EThreeKSSixRover}{11.13}
\newcommand{\EThreeKSSixRoverIrr}{14.6}
\newcommand{\EThreeKSTwoAtt}{64}
\newcommand{\EThreeKSTwoCo}{0.17}
\newcommand{\EThreeKSTwoCoHi}{0.74}
\newcommand{\EThreeKSTwoCoIrr}{-0.8}
\newcommand{\EThreeKSTwoCoIrrHi}{1.2}
\newcommand{\EThreeKSTwoCoIrrLo}{-3.8}
\newcommand{\EThreeKSTwoCoIrrP}{=0.458}
\newcommand{\EThreeKSTwoCoLo}{-0.53}
\newcommand{\EThreeKSTwoCoLoses}{7}
\newcommand{\EThreeKSTwoCoP}{=0.012}
\newcommand{\EThreeKSTwoCoWins}{19}
\newcommand{\EThreeKSTwoNa}{0.68}
\newcommand{\EThreeKSTwoNaHi}{1.20}
\newcommand{\EThreeKSTwoNaIrr}{4.6}
\newcommand{\EThreeKSTwoNaIrrHi}{7.1}
\newcommand{\EThreeKSTwoNaIrrLo}{2.1}
\newcommand{\EThreeKSTwoNaIrrP}{=0.002}
\newcommand{\EThreeKSTwoNaLo}{0.16}
\newcommand{\EThreeKSTwoNaLoses}{12}
\newcommand{\EThreeKSTwoNaP}{=0.033}
\newcommand{\EThreeKSTwoNaWins}{17}
\newcommand{\EThreeKSTwoNo}{1.65}
\newcommand{\EThreeKSTwoNoHi}{2.43}
\newcommand{\EThreeKSTwoNoIrr}{10.8}
\newcommand{\EThreeKSTwoNoIrrHi}{14.2}
\newcommand{\EThreeKSTwoNoIrrLo}{7.9}
\newcommand{\EThreeKSTwoNoIrrP}{<0.001}
\newcommand{\EThreeKSTwoNoLo}{0.90}
\newcommand{\EThreeKSTwoNoLoses}{9}
\newcommand{\EThreeKSTwoNoP}{<0.001}
\newcommand{\EThreeKSTwoNoWins}{21}
\newcommand{\EThreeKSTwoRb}{1.66}
\newcommand{\EThreeKSTwoRbHi}{2.55}
\newcommand{\EThreeKSTwoRbIrr}{2.9}
\newcommand{\EThreeKSTwoRbIrrHi}{6.2}
\newcommand{\EThreeKSTwoRbIrrLo}{-0.4}
\newcommand{\EThreeKSTwoRbIrrP}{=0.112}
\newcommand{\EThreeKSTwoRbLo}{0.80}
\newcommand{\EThreeKSTwoRbLoses}{9}
\newcommand{\EThreeKSTwoRbP}{=0.002}
\newcommand{\EThreeKSTwoRbWins}{20}
\newcommand{\EThreeKSTwoRover}{10.26}
\newcommand{\EThreeKSTwoRoverIrr}{12.1}
\newcommand{\EThreeKSZeroAtt}{64}
\newcommand{\EThreeKSZeroCo}{0.58}
\newcommand{\EThreeKSZeroCoHi}{0.98}
\newcommand{\EThreeKSZeroCoIrr}{0.8}
\newcommand{\EThreeKSZeroCoIrrHi}{2.1}
\newcommand{\EThreeKSZeroCoIrrLo}{0.0}
\newcommand{\EThreeKSZeroCoIrrP}{=0.157}
\newcommand{\EThreeKSZeroCoLo}{0.23}
\newcommand{\EThreeKSZeroCoLoses}{4}
\newcommand{\EThreeKSZeroCoP}{<0.001}
\newcommand{\EThreeKSZeroCoWins}{19}
\newcommand{\EThreeKSZeroMy}{2.22}
\newcommand{\EThreeKSZeroMyHi}{3.31}
\newcommand{\EThreeKSZeroMyIrr}{10.0}
\newcommand{\EThreeKSZeroMyIrrHi}{13.8}
\newcommand{\EThreeKSZeroMyIrrLo}{6.7}
\newcommand{\EThreeKSZeroMyIrrP}{<0.001}
\newcommand{\EThreeKSZeroMyLo}{1.24}
\newcommand{\EThreeKSZeroMyLoses}{8}
\newcommand{\EThreeKSZeroMyP}{<0.001}
\newcommand{\EThreeKSZeroMyWins}{22}
\newcommand{\EThreeKSZeroNa}{1.09}
\newcommand{\EThreeKSZeroNaHi}{1.85}
\newcommand{\EThreeKSZeroNaIrr}{5.4}
\newcommand{\EThreeKSZeroNaIrrHi}{8.3}
\newcommand{\EThreeKSZeroNaIrrLo}{2.9}
\newcommand{\EThreeKSZeroNaIrrP}{=0.002}
\newcommand{\EThreeKSZeroNaLo}{0.39}
\newcommand{\EThreeKSZeroNaLoses}{10}
\newcommand{\EThreeKSZeroNaP}{=0.008}
\newcommand{\EThreeKSZeroNaWins}{18}
\newcommand{\EThreeKSZeroNo}{2.58}
\newcommand{\EThreeKSZeroNoHi}{3.80}
\newcommand{\EThreeKSZeroNoIrr}{13.3}
\newcommand{\EThreeKSZeroNoIrrHi}{17.5}
\newcommand{\EThreeKSZeroNoIrrLo}{9.6}
\newcommand{\EThreeKSZeroNoIrrP}{<0.001}
\newcommand{\EThreeKSZeroNoLo}{1.60}
\newcommand{\EThreeKSZeroNoLoses}{5}
\newcommand{\EThreeKSZeroNoP}{<0.001}
\newcommand{\EThreeKSZeroNoWins}{25}
\newcommand{\EThreeKSZeroRb}{2.08}
\newcommand{\EThreeKSZeroRbHi}{2.95}
\newcommand{\EThreeKSZeroRbIrr}{4.6}
\newcommand{\EThreeKSZeroRbIrrHi}{7.1}
\newcommand{\EThreeKSZeroRbIrrLo}{2.1}
\newcommand{\EThreeKSZeroRbIrrP}{=0.002}
\newcommand{\EThreeKSZeroRbLo}{1.27}
\newcommand{\EThreeKSZeroRbLoses}{4}
\newcommand{\EThreeKSZeroRbP}{<0.001}
\newcommand{\EThreeKSZeroRbWins}{26}
\newcommand{\EThreeKSZeroRover}{9.28}
\newcommand{\EThreeKSZeroRoverIrr}{10.4}
\newcommand{\EThreeLBl}{11.18}
\newcommand{\EThreeLCo}{9.86}
\newcommand{\EThreeLMy}{11.50}
\newcommand{\EThreeLNa}{10.37}
\newcommand{\EThreeLNo}{11.86}
\newcommand{\EThreeLOneBlind}{11.18}
\newcommand{\EThreeLOneConst}{9.86}
\newcommand{\EThreeLOneMyopic}{11.50}
\newcommand{\EThreeLOneNoAudit}{11.86}
\newcommand{\EThreeLOneNoaudit}{10.37}
\newcommand{\EThreeLOnePeriodic}{16.07}
\newcommand{\EThreeLOneRandom}{13.46}
\newcommand{\EThreeLOneReset}{25.94}
\newcommand{\EThreeLOneRollback}{11.36}
\newcommand{\EThreeLOneRover}{9.28}
\newcommand{\EThreeLOneValidator}{15.19}
\newcommand{\EThreeLRb}{11.36}
\newcommand{\EThreeLTenBlind}{19.99}
\newcommand{\EThreeLTenConst}{18.79}
\newcommand{\EThreeLTenMyopic}{22.38}
\newcommand{\EThreeLTenNoAudit}{20.69}
\newcommand{\EThreeLTenNoaudit}{19.94}
\newcommand{\EThreeLTenPeriodic}{26.06}
\newcommand{\EThreeLTenRandom}{23.36}
\newcommand{\EThreeLTenReplayOnly}{39.64}
\newcommand{\EThreeLTenReset}{25.94}
\newcommand{\EThreeLTenRollback}{23.02}
\newcommand{\EThreeLTenRover}{18.82}
\newcommand{\EThreeLTenValidator}{22.77}
\newcommand{\EThreeLTenthBlind}{6.63}
\newcommand{\EThreeLTenthConst}{6.17}
\newcommand{\EThreeLTenthMyopic}{7.76}
\newcommand{\EThreeLTenthNoAudit}{8.10}
\newcommand{\EThreeLTenthNoaudit}{6.09}
\newcommand{\EThreeLTenthPeriodic}{12.06}
\newcommand{\EThreeLTenthRandom}{8.73}
\newcommand{\EThreeLTenthReset}{25.94}
\newcommand{\EThreeLTenthRollback}{5.94}
\newcommand{\EThreeLTenthRover}{5.40}
\newcommand{\EThreeLTenthValidator}{11.18}
\newcommand{\EThreeLZeroBlind}{6.49}
\newcommand{\EThreeLZeroConst}{5.90}
\newcommand{\EThreeLZeroMyopic}{7.34}
\newcommand{\EThreeLZeroNoAudit}{7.68}
\newcommand{\EThreeLZeroNoaudit}{5.41}
\newcommand{\EThreeLZeroPeriodic}{11.61}
\newcommand{\EThreeLZeroRandom}{8.14}
\newcommand{\EThreeLZeroReset}{25.94}
\newcommand{\EThreeLZeroRollback}{5.33}
\newcommand{\EThreeLZeroRover}{4.85}
\newcommand{\EThreeLZeroValidator}{10.73}
\newcommand{\EThreeMaxBaseP}{0.002}
\newcommand{\EThreeMsPerRecovery}{1542}
\newcommand{\EThreeNTests}{10}
\newcommand{\EThreePOneBlind}{=0.002}
\newcommand{\EThreePOneConst}{=0.002}
\newcommand{\EThreePOneMyopic}{=0.002}
\newcommand{\EThreePOneNoAudit}{<0.001}
\newcommand{\EThreePOneNoaudit}{=0.008}
\newcommand{\EThreePOnePeriodic}{<0.001}
\newcommand{\EThreePOneRandom}{<0.001}
\newcommand{\EThreePOneReset}{<0.001}
\newcommand{\EThreePOneRollback}{<0.001}
\newcommand{\EThreePOneValidator}{<0.001}
\newcommand{\EThreePTenBlind}{=0.229}
\newcommand{\EThreePTenConst}{=0.927}
\newcommand{\EThreePTenMyopic}{<0.001}
\newcommand{\EThreePTenNoAudit}{=0.126}
\newcommand{\EThreePTenNoaudit}{=0.159}
\newcommand{\EThreePTenPeriodic}{<0.001}
\newcommand{\EThreePTenRandom}{=0.002}
\newcommand{\EThreePTenReplayOnly}{<0.001}
\newcommand{\EThreePTenReset}{<0.001}
\newcommand{\EThreePTenRollback}{<0.001}
\newcommand{\EThreePTenValidator}{=0.002}
\newcommand{\EThreePTenthBlind}{=0.030}
\newcommand{\EThreePTenthConst}{<0.001}
\newcommand{\EThreePTenthMyopic}{<0.001}
\newcommand{\EThreePTenthNoAudit}{<0.001}
\newcommand{\EThreePTenthNoaudit}{=0.238}
\newcommand{\EThreePTenthPeriodic}{<0.001}
\newcommand{\EThreePTenthRandom}{<0.001}
\newcommand{\EThreePTenthReset}{<0.001}
\newcommand{\EThreePTenthRollback}{=0.328}
\newcommand{\EThreePTenthValidator}{<0.001}
\newcommand{\EThreePZeroBlind}{=0.002}
\newcommand{\EThreePZeroConst}{=0.001}
\newcommand{\EThreePZeroMyopic}{=0.001}
\newcommand{\EThreePZeroNoAudit}{<0.001}
\newcommand{\EThreePZeroNoaudit}{=0.300}
\newcommand{\EThreePZeroPeriodic}{<0.001}
\newcommand{\EThreePZeroRandom}{<0.001}
\newcommand{\EThreePZeroReset}{<0.001}
\newcommand{\EThreePZeroRollback}{=0.360}
\newcommand{\EThreePZeroValidator}{<0.001}
\newcommand{\EThreePerCell}{240}
\newcommand{\EThreePrimHOneAtw}{0.37}
\newcommand{\EThreePrimHOneP}{<0.001}
\newcommand{\EThreePrimHTwoAtw}{0.40}
\newcommand{\EThreePrimHTwoP}{=0.002}
\newcommand{\EThreeRbCheapTenth}{6.06}
\newcommand{\EThreeRbCheapTenthD}{0.66}
\newcommand{\EThreeRbCheapTenthHi}{1.15}
\newcommand{\EThreeRbCheapTenthLo}{0.17}
\newcommand{\EThreeRbCheapTenthP}{=0.109}
\newcommand{\EThreeRbCheapZero}{5.38}
\newcommand{\EThreeRbCheapZeroD}{0.53}
\newcommand{\EThreeRbCheapZeroHi}{1.03}
\newcommand{\EThreeRbCheapZeroLo}{0.02}
\newcommand{\EThreeRbCheapZeroP}{=0.338}
\newcommand{\EThreeRbFr}{0.083}
\newcommand{\EThreeRbFrD}{0.050}
\newcommand{\EThreeRbFrHi}{0.096}
\newcommand{\EThreeRbFrLo}{0.008}
\newcommand{\EThreeRbFrP}{=0.042}
\newcommand{\EThreeRbFrRover}{0.033}
\newcommand{\EThreeRbRec}{0.69}
\newcommand{\EThreeRbRecD}{0.07}
\newcommand{\EThreeRbRecHi}{0.12}
\newcommand{\EThreeRbRecLo}{0.03}
\newcommand{\EThreeRbRecP}{=0.005}
\newcommand{\EThreeRbRecRover}{0.62}
\newcommand{\EThreeRbRegen}{6.03}
\newcommand{\EThreeRbRegenD}{2.09}
\newcommand{\EThreeRbRegenHi}{2.70}
\newcommand{\EThreeRbRegenLo}{1.53}
\newcommand{\EThreeRbRegenP}{<0.001}
\newcommand{\EThreeRbRegenPer}{8.7}
\newcommand{\EThreeRbRegenRover}{3.94}
\newcommand{\EThreeRecoveries}{7{,}935}
\newcommand{\EThreeRegen}{54{,}298}
\newcommand{\EThreeRemoteCanary}{10{,}016}
\newcommand{\EThreeResetCleanPct}{56}
\newcommand{\EThreeRobDalBest}{11.36}
\newcommand{\EThreeRobDalBestName}{rollback}
\newcommand{\EThreeRobDalD}{-0.89}
\newcommand{\EThreeRobDalDhi}{-0.08}
\newcommand{\EThreeRobDalDlo}{-1.74}
\newcommand{\EThreeRobDalMyopic}{11.57}
\newcommand{\EThreeRobDalNoAudit}{11.86}
\newcommand{\EThreeRobDalRollback}{11.36}
\newcommand{\EThreeRobDalRover}{10.47}
\newcommand{\EThreeRobLateBest}{10.06}
\newcommand{\EThreeRobLateBestName}{no audit}
\newcommand{\EThreeRobLateConst}{10.27}
\newcommand{\EThreeRobLateD}{0.53}
\newcommand{\EThreeRobLateDhi}{1.13}
\newcommand{\EThreeRobLateDlo}{-0.10}
\newcommand{\EThreeRobLateMyopic}{10.13}
\newcommand{\EThreeRobLateNoAudit}{10.06}
\newcommand{\EThreeRobLateNoaudit}{10.19}
\newcommand{\EThreeRobLateRollback}{12.04}
\newcommand{\EThreeRobLateRover}{10.59}
\newcommand{\EThreeRobSoneBest}{11.76}
\newcommand{\EThreeRobSoneBestName}{myopic}
\newcommand{\EThreeRobSoneD}{-0.86}
\newcommand{\EThreeRobSoneDhi}{0.05}
\newcommand{\EThreeRobSoneDlo}{-1.88}
\newcommand{\EThreeRobSoneMyopic}{11.76}
\newcommand{\EThreeRobSoneNoAudit}{11.86}
\newcommand{\EThreeRobSoneRollback}{12.50}
\newcommand{\EThreeRobSoneRover}{10.90}
\newcommand{\EThreeRoverAudits}{0.90}
\newcommand{\EThreeRoverHi}{10.39}
\newcommand{\EThreeRoverLo}{8.26}
\newcommand{\EThreeRoverLoss}{9.28}
\newcommand{\EThreeRoverRegenPer}{6.4}
\newcommand{\EThreeSeeds}{8}
\newcommand{\EThreeShareTenBranch}{29}
\newcommand{\EThreeShareTenMemory}{17}
\newcommand{\EThreeShareTenQueue}{24}
\newcommand{\EThreeShareTenSkill}{30}
\newcommand{\EThreeShareZeroBranch}{7}
\newcommand{\EThreeShareZeroMemory}{9}
\newcommand{\EThreeShareZeroQueue}{28}
\newcommand{\EThreeShareZeroSkill}{56}
\newcommand{\EThreeStreams}{30}
\newcommand{\EThreeSwTenBest}{20.69}
\newcommand{\EThreeSwTenBestD}{1.87}
\newcommand{\EThreeSwTenBestDhi}{3.62}
\newcommand{\EThreeSwTenBestDlo}{0.19}
\newcommand{\EThreeSwTenBestName}{no audit}
\newcommand{\EThreeSwTenBestP}{=0.126}
\newcommand{\EThreeSwTenthBest}{5.94}
\newcommand{\EThreeSwTenthBestD}{0.53}
\newcommand{\EThreeSwTenthBestDhi}{1.19}
\newcommand{\EThreeSwTenthBestDlo}{-0.07}
\newcommand{\EThreeSwTenthBestName}{rollback only}
\newcommand{\EThreeSwTenthBestP}{=0.328}
\newcommand{\EThreeSwZeroBest}{5.33}
\newcommand{\EThreeSwZeroBestD}{0.48}
\newcommand{\EThreeSwZeroBestDhi}{1.19}
\newcommand{\EThreeSwZeroBestDlo}{-0.17}
\newcommand{\EThreeSwZeroBestName}{rollback only}
\newcommand{\EThreeSwZeroBestP}{=0.360}
\newcommand{\EThreeSymbolic}{No audit,Random,Periodic,Rollback,Reset all,Validator,Myopic,ROVER phi-blind,ROVER no window,ROVER no audits,ROVER}
\newcommand{\EThreeThrCheap}{0.5}
\newcommand{\EThreeThrOne}{0.7}
\newcommand{\EThreeThrTen}{0.7}
\newcommand{\EThreeTuneEpisodes}{3{,}360}
\newcommand{\EThreeVOneResetTen}{28.39}
\newcommand{\EThreeVOneRollbackZero}{5.28}
\newcommand{\EThreeVOneRoverTen}{42.43}
\newcommand{\EThreeVOneRoverZero}{5.70}
\newcommand{\EThreeVTwoResetTen}{23.34}
\newcommand{\EThreeVTwoRollbackZero}{4.81}
\newcommand{\EThreeVTwoRoverTen}{17.72}
\newcommand{\EThreeVTwoRoverZero}{4.64}
\newcommand{\EThreeVsBestAtw}{0.63}
\newcommand{\EThreeVsBestD}{2.08}
\newcommand{\EThreeVsBestDhi}{2.95}
\newcommand{\EThreeVsBestDlo}{1.27}
\newcommand{\EThreeVsBestP}{<0.001}
\newcommand{\EThreeVsBestPct}{18.3}
\newcommand{\EThreeVsBestWins}{26}
\newcommand{\ETwoBRBlindOne}{queue}
\newcommand{\ETwoBRBlindTen}{branch}
\newcommand{\ETwoBRBlindZero}{queue}
\newcommand{\ETwoBRPhiZero}{memory}
\newcommand{\ETwoBlindRatio}{3.3}
\newcommand{\ETwoMCcells}{16}
\newcommand{\ETwoMCinside}{15}
\newcommand{\ETwoMCmaxdev}{0.286}
\newcommand{\ETwoQueueEnterHi}{10}
\newcommand{\ETwoQueueEnterLo}{7}
\newcommand{\ETwoUAwareOne}{5.97}
\newcommand{\ETwoUAwareTen}{22.71}
\newcommand{\ETwoUAwareZero}{4.11}
\newcommand{\ETwoUBlindOne}{14.27}
\newcommand{\ETwoUBlindZero}{13.76}
\newcommand{\ETwoUPhiTen}{116.70}
\newcommand{\ETwoUPhiZero}{65.76}
\newcommand{\ETwoUUniformZero}{5.70}
\newcommand{\ETwoUstarTen}{22.71}
\newcommand{\ETwoUstarZero}{4.11}
\newcommand{\ETwoXBranchOne}{0.143}
\newcommand{\ETwoXBranchTen}{0.161}
\newcommand{\ETwoXBranchZero}{0.138}
\newcommand{\ETwoXMemoryOne}{0.218}
\newcommand{\ETwoXMemoryTen}{0.311}
\newcommand{\ETwoXMemoryZero}{0.179}
\newcommand{\ETwoXQueueOne}{0.072}
\newcommand{\ETwoXQueueTen}{0.000}
\newcommand{\ETwoXQueueZero}{0.096}
\newcommand{\ETwoXSkillOne}{0.167}
\newcommand{\ETwoXSkillTen}{0.128}
\newcommand{\ETwoXSkillZero}{0.187}
\newcommand{\GridMaxRel}{$5\times10^{-5}$}
\newcommand{\GridMaxT}{6}
\newcommand{\OvDecideUs}{292}
\newcommand{\OvSolveSec}{0.08}
\newcommand{\PThreeAdm}{1{,}187}
\newcommand{\PThreeAdmLoses}{392}
\newcommand{\PThreeExRatio}{10.7}
\newcommand{\PThreeExSBaZero}{1.66}
\newcommand{\PThreeExSBbZero}{0.16}
\newcommand{\PThreeExXaTwo}{0.48}
\newcommand{\PThreeExXaZero}{0.24}
\newcommand{\PThreeInstances}{1{,}497}
\newcommand{\PThreeKinkLeft}{$-0.0076$}
\newcommand{\PThreeKinkR}{7.27}
\newcommand{\PThreeKinkRight}{$0.0$}
\newcommand{\PThreeLowHbr}{skill}
\newcommand{\PThreeLowHbrU}{9.24}
\newcommand{\PThreeLowHleast}{queued call}
\newcommand{\PThreeLowHleastU}{5.80}
\newcommand{\PThreeNotArgmax}{784}
\newcommand{\PThreePhiLoses}{493}
\newcommand{\PThreeUmemory}{65.76}
\newcommand{\PThreeUqueue}{13.76}
\newcommand{\PThreeUskill}{18.44}
\newcommand{\PTwoAuditInst}{16}
\newcommand{\PTwoConfirmed}{41}
\newcommand{\PTwoPatternInst}{39}
\newcommand{\PTwoRecheckAudit}{16}
\newcommand{\PTwoRecheckPattern}{25}
\newcommand{\PTwoThrAInst}{571}
\newcommand{\PTwoThrRInst}{556}
\newcommand{\PTwoTried}{41}
\newcommand{\PTwoUpperInst}{0}
\newcommand{\PostATenOK}{yes}
\newcommand{\PostAdNMax}{45}
\newcommand{\PostAdNMin}{30}
\newcommand{\PostAdRoSplit}{8.5}
\newcommand{\PostAdVaSplit}{9.8}
\newcommand{\PostCsConfMax}{0.21}
\newcommand{\PostCsConfMin}{0.12}
\newcommand{\PostCsCostErrMax}{1.8}
\newcommand{\PostCsDirectMax}{1.09}
\newcommand{\PostCsDirectMin}{0.94}
\newcommand{\PostCsEpisodes}{960}
\newcommand{\PostCsGrowMax}{22}
\newcommand{\PostCsGrowMin}{8}
\newcommand{\PostCsResLowMax}{28}
\newcommand{\PostCsResLowMin}{12}
\newcommand{\PostGridActDis}{0.002}
\newcommand{\PostGridDense}{2{,}001}
\newcommand{\PostGridEps}{960}
\newcommand{\PostGridIdent}{960}
\newcommand{\PostGridMaxErr}{0.0028}
\newcommand{\PostGridMaxRel}{0.15}
\newcommand{\PostGridRef}{8{,}001}
\newcommand{\PostOsBlindHelps}{18}
\newcommand{\PostOsBlindHurts}{25}
\newcommand{\PostOsCells}{108}
\newcommand{\PostOsMcZ}{2.37}
\newcommand{\PostOsMyGap}{69.6}
\newcommand{\PostOsRollBeats}{101}
\newcommand{\PostOsRollGap}{42.5}
\newcommand{\PostOsRollHighMax}{4.9}
\newcommand{\PostOsTermGap}{61.8}
\newcommand{\PostPubEpisodes}{4{,}320}
\newcommand{\PostPubFreqLineFivePe}{17.75}
\newcommand{\PostPubFreqLineFiveRank}{1}
\newcommand{\PostPubFreqLineFiveRo}{14.38}
\newcommand{\PostPubFreqLineFiveVa}{15.38}
\newcommand{\PostPubFreqLineFiveVaD}{1.00}
\newcommand{\PostPubFreqLineFiveVaHi}{2.49}
\newcommand{\PostPubFreqLineFiveVaLo}{-0.74}
\newcommand{\PostPubFreqLineTwPe}{22.82}
\newcommand{\PostPubFreqLineTwRank}{3}
\newcommand{\PostPubFreqLineTwRo}{29.70}
\newcommand{\PostPubFreqLineTwVa}{15.94}
\newcommand{\PostPubFreqLineTwVaD}{-13.75}
\newcommand{\PostPubFreqLineTwVaHi}{-8.58}
\newcommand{\PostPubFreqLineTwVaLo}{-19.54}
\newcommand{\PostPubFreqNoLeak}{46}
\newcommand{\PostPubFreqNoLines}{2.00}
\newcommand{\PostPubFreqRelFivePe}{17.82}
\newcommand{\PostPubFreqRelFiveRank}{1}
\newcommand{\PostPubFreqRelFiveRo}{13.11}
\newcommand{\PostPubFreqRelFiveVa}{15.61}
\newcommand{\PostPubFreqRelFiveVaD}{2.50}
\newcommand{\PostPubFreqRelFiveVaHi}{3.92}
\newcommand{\PostPubFreqRelFiveVaLo}{0.80}
\newcommand{\PostPubFreqRelTwPe}{23.07}
\newcommand{\PostPubFreqRelTwRank}{3}
\newcommand{\PostPubFreqRelTwRo}{24.61}
\newcommand{\PostPubFreqRelTwVa}{16.86}
\newcommand{\PostPubFreqRelTwVaD}{-7.75}
\newcommand{\PostPubFreqRelTwVaHi}{-3.38}
\newcommand{\PostPubFreqRelTwVaLo}{-12.93}
\newcommand{\PostPubFreqRoLeak}{30}
\newcommand{\PostPubFreqRoLines}{1.02}
\newcommand{\PostPubFreqVaLeak}{5}
\newcommand{\PostPubFreqVaLines}{0.04}
\newcommand{\PostPubFreqWorstRank}{3}
\newcommand{\PostPubRareAllFirst}{yes}
\newcommand{\PostPubRareNoLeak}{20}
\newcommand{\PostPubRareNoLines}{0.70}
\newcommand{\PostPubRareRoLeak}{8}
\newcommand{\PostPubRareRoLines}{0.26}
\newcommand{\PostPubRareVaLeak}{2}
\newcommand{\PostPubRareVaLines}{0.02}
\newcommand{\PostPubRareWorstRank}{1}
\newcommand{\PostPubRatioMax}{3.1}
\newcommand{\PostPubRatioRo}{2.7}
\newcommand{\PostStFlips}{1}
\newcommand{\PostStHOneAtw}{0.63}
\newcommand{\PostStHOneSf}{<0.001}
\newcommand{\PostStHOneT}{<0.001}
\newcommand{\PostStHOneWin}{0.87}
\newcommand{\PostStHTwoSf}{<0.001}
\newcommand{\PostStHTwoT}{<0.001}
\newcommand{\PostStTests}{41}
\newcommand{\PropAParams}{$\lambda=0.2$, $\phi=0$, $H=10$, $h=4$, $r_0=4$, $r_1=1$, $c_a=2$, $\sigma_0=0.2$, $\sigma_1=0.9$, $D_0=10$}
\newcommand{\PropBBelief}{0.4}
\newcommand{\PropBNpol}{63{,}756}
\newcommand{\PropBParams}{$\lambda=0.3$, $\phi=1$, $H=20$, $h=1$, $r_0=1$, $r_1=1$, $c_a=2$, $d=0.9$, $\alpha=0.05$, $\sigma_0=0.1$, $\sigma_1=0.4$}
\newcommand{\PropBThrSix}{0.597}
\newcommand{\PropBThrSixExact}{0.5968}
\newcommand{\PropBThrTwo}{0.213}
\newcommand{\PropBThrTwoExact}{0.2130}
\newcommand{\PropBa}{0.20}
\newcommand{\PropBb}{0.55}
\newcommand{\PropBc}{0.82}
\newcommand{\PropBd}{0.95}
\newcommand{\PropCBeliefs}{0.05, 0.18, 0.33, 0.6, 0.95}
\newcommand{\PropCParams}{$\lambda=0.2$, $\phi=0.5$, $H=2$, $h=4$, $r_0=2$, $r_1=2$, $c_a=0.25$, $d=0.8$, $\alpha=0.3$, $\sigma_0=0.2$, $\sigma_1=0.7$, $D_0=6$}
\newcommand{\PropCa}{2}
\newcommand{\PropDzero}{10}
\newcommand{\PropH}{4}
\newcommand{\PropLam}{0.2}
\newcommand{\PropNpol}{18{,}396}
\newcommand{\PropPostHi}{0.94}
\newcommand{\PropPostLo}{0.31}
\newcommand{\PropPostMid}{0.81}
\newcommand{\PropPostMidExact}{0.8148}
\newcommand{\PropRone}{1}
\newcommand{\PropRzero}{4}
\newcommand{\PropSegAend}{0.755}
\newcommand{\PropSegAstart}{0.385}
\newcommand{\PropSegWend}{0.383}
\newcommand{\PropSegWtwoend}{0.890}
\newcommand{\PropSegWtwostart}{0.758}
\newcommand{\PropSone}{0.9}
\newcommand{\PropSzero}{0.2}
\newcommand{\TOneAsixInstances}{234}
\newcommand{\TOneCells}{36{,}247}
\newcommand{\TOneInstances}{1{,}000}
\newcommand{\TOneIntervalViol}{0}
\newcommand{\TOneMonoCells}{7{,}166}
\newcommand{\TOneMonoViol}{0}
\newcommand{\TOnePAInstances}{400}
\newcommand{\TOnePAViol}{0}
\newcommand{\TThreeAfail}{40}
\newcommand{\TThreeBRn}{1{,}500}
\newcommand{\TThreeBRviol}{0}
\newcommand{\TThreeDmaxrel}{$10^{-6}$}
\newcommand{\TThreeDn}{352}
\newcommand{\TThreeDok}{352}
\newcommand{\TThreeDsign}{352}
\newcommand{\TThreeMCn}{80}
\newcommand{\TThreeMCok}{80}
\newcommand{\TThreeNonpos}{8}
\newcommand{\TTwoBoundViol}{0}
\newcommand{\TTwoCornerViol}{0}
\newcommand{\TTwoDefers}{339}
\newcommand{\TTwoGapOverBound}{0.95}
\newcommand{\TTwoInstances}{600}
\newcommand{\TTwoMaxErr}{$3\times10^{-15}$}
\newcommand{\TTwoMyMatch}{600}
\newcommand{\TTwoRatioEighty}{89.0}
\newcommand{\TTwoRatioFive}{10.4}
\newcommand{\TTwoTightRatio}{0.950}
\newcommand{\ThrAtenDten}{0.27}
\newcommand{\ThrAtenDzero}{0.22}
\newcommand{\ThrRtenDsix}{0.81}
\newcommand{\ThrRtenDten}{0.86}
\newcommand{\ThrRtenDzero}{0.52}


\newtheorem{theorem}{Theorem}
\newtheorem{proposition}[theorem]{Proposition}
\newtheorem{lemma}{Lemma}
\theoremstyle{definition}
\newtheorem{assumption}{Assumption}

\newcommand{\Rc}{\mathcal{R}}
\newcommand{\Ac}{\mathcal{A}}
\newcommand{\E}{\mathbb{E}}
\renewcommand{\P}{\mathbb{P}}
\newcommand{\Wa}{\mathsf{W}}
\newcommand{\Aa}{\mathsf{A}}
\newcommand{\Ra}{\mathsf{R}}

\title{The Option Value of Rollback\\[2pt]\large Supplementary material: full proofs and verification protocol}
\author{Anonymous submission}
\date{}

\begin{document}
\maketitle

\noindent This supplement contains the complete proofs of the formal results stated in the paper,
the exact statement of every assumption, and the protocol of the numerical checks that accompany
each result. Numbering of theorems, propositions, definitions and assumptions matches the paper.
Every number quoted here is generated by the artifact.

\section{Model, notation and assumptions}

Periods are $t=0,\dots,T-1$. A single carrier has a latent taint bit $z_t\in\{0,1\}$ and an observed
amount of derived work $D_t\in\mathbb{N}$ accumulated since the last verified checkpoint. Recovery
costs $R(D)=r_0+r_1D$. In each period the defender chooses wait ($\Wa$), recover ($\Ra$: roll back to
the checkpoint and replay clean work), or a contingent audit ($\Aa$: pay $c_a$, observe the outcome,
then choose $\Wa$ or $\Ra$ within the same period). Period $t$ runs in five steps: (1) action, (2) growth of derived work by $\xi_t$, (3) activation of a
tainted carrier that was not recovered, with probability $\lambda$, (4) otherwise a monitor output and
a belief update, and (5) infection of a clean carrier with probability $\rho_t$.

\begin{assumption}[A1, taint]
$\P(z_0=1)=\pi$. At the end of period $t$ a clean carrier becomes tainted with probability
$\rho_t\in[0,1)$. Taint persists until a recovery, which sets $z=0$.
\end{assumption}
\begin{assumption}[A2, derived work]
$D_{t+1}=D_t+\xi_t$ unless a recovery occurs in period $t$, where the increments are i.i.d.,
nonnegative, integer valued with finite support and mean $\Delta$, and independent of all other
randomness. A chosen recovery in period $t$ resets $D$ before the period's growth, which gives
$D_{t+1}=\xi_t$. A forced recovery after an activation (A4) resets it after the growth, which gives
$D_{t+1}=0$. No result depends on this convention.
\end{assumption}
\begin{assumption}[A3, sensing]
The monitor emits $o_t\in\{0,1\}$ with $\P(o_t=1\mid z=1)=\sigma_1$ and $\P(o_t=1\mid z=0)=\sigma_0$.
An audit is positive with probability $d$ if $z=1$ and $\alpha$ if $z=0$. Given $z$ all observations
are independent. The monotone likelihood ratio $\sigma_1\ge\sigma_0$ motivates the conjecture of
Definition 3 but is not used by any proof.
\end{assumption}
\begin{assumption}[A4, activation]
A tainted carrier that is not recovered in period $t$ activates with probability $\lambda\in(0,1)$.
Activation is observed, costs $\phi H$ (irreversible harm, $\phi\in[0,1]$, $H\ge0$) and triggers a
mandatory recovery at cost $R(D_t+\xi_t)$, after which $z=0$ and $D=0$.
\end{assumption}
\begin{assumption}[A5, costs and horizon]
$T<\infty$, all cost parameters are nonnegative, recovering a clean carrier costs an extra $c_{fr}$, and
the terminal cost is $hD_Tz_T$. The objective is the expected total cost.
\end{assumption}
\begin{assumption}[A6] $h\le r_1$. \end{assumption}
\begin{assumption}[A7]
$d=1$, $\alpha=0$, $\rho_t\equiv0$, $\sigma_0=\sigma_1$, $c_{fr}=0$, $\xi_t\equiv\Delta>0$, and
$R(D)\le hD$ for all $D\ge D_0$. A6 and A7 hold together only if $h=r_1$ and $r_0=0$.
\end{assumption}
\begin{assumption}[A8, commitment game]
The defender commits to per-period audit probabilities $x\in X_B=\{x\in[0,1]^n:\sum_ix_i\le B\}$. The
attacker observes $x$ and poisons one carrier $i$ at $t=0$, succeeding with probability $s_i\in(0,1]$ at
cost $\kappa_i\ge0$, or abstains with payoff $0$. Its payoff is $U_i(x_i)=s_iL_i(x_i,r_1)-\kappa_i$, where
$L_i$ is the defender's expected cost caused by a successful poison.
\end{assumption}
\begin{assumption}[A9]
Carrier $i$ has $D=0$ when it is poisoned at $t=0$, grows derived work deterministically by $\Delta_i$
per period and follows A4 and A5 with $(\lambda_i,\phi_i)$. In each period it is audited independently with probability $x_i$. An audit of a
tainted carrier detects it with probability $d$ and triggers recovery at cost $R(D_t)$. An undetected
taint at $T$ leaves residual harm $hD_T$. Moreover $h\ge r_1$, under which leaving a detected taint in
place never saves more than $r_0$ (Section 4).
\end{assumption}
\begin{assumption}[A10, regularity at $r_1$]
$\partial U_i/\partial x_i<0$ on $[0,\min\{1,B\}]$ for every $i$. At a minimax commitment $x^*$ with
level $u^*=\max_iU_i(x^*_i)$ and attack set $S=\{i:U_i(x_i^*)=u^*\}$: $u^*>0$, $0<x^*_i<1$ for $i\in S$,
$\sum_ix_i^*=B$, and $U_j(0)<u^*$ for $j\notin S$. The condition $u^*>0$ means that the attacker strictly
prefers to attack. Without it every commitment with $\max_iU_i\le0$ attains the value $0$ of
Definition 5, and the minimax commitment need not be unique.
\end{assumption}

\paragraph{Belief update.} With $b=\P(z=1\mid\text{history})$, the posterior after no activation and
monitor output $o$ is $b^o=b(1-\lambda)\ell_1(o)/[b(1-\lambda)\ell_1(o)+(1-b)\ell_0(o)]$ with
$\ell_z(1)=\sigma_z$ and $\ell_z(0)=1-\sigma_z$, followed by $b\mapsto b+(1-b)\rho_t$. An audit outcome
updates $b$ to $b^+=bd/(bd+(1-b)\alpha)$ or $b^-=b(1-d)/(b(1-d)+(1-b)(1-\alpha))$.

\paragraph{Plans.} From a decision point $(t,D,b)$, a deterministic \emph{plan} $\sigma$ maps every finite
sequence of future observations (monitor outputs, audit outcomes, activation indicators and work
increments) to an action. Under A2 and A3 the set $\Sigma_t$ of plans is finite and does not depend on
$(D,b)$. Write $J_\sigma(t,D,b)$ for its expected cost, $Q_a(t,D,b)$ for the minimum over plans whose first
action is $a$, and $V_t=\min_aQ_a$. Optimal history-dependent policies are attained by plans because, for a
fixed initial condition, $D_s$ and $b_s$ are functions of the observation history. The recovery region is
$\Rc_t(D)=\{b:Q_\Ra\le\min(Q_\Wa,Q_\Aa)\}$ and the audit region $\Ac_t(D)$ is the set where $\Aa$ attains
the minimum.

\section{Structure of the optimal adaptation policy}

\begin{lemma}[affine decomposition]\label{lem:affine}
Assume A1 to A5. For every plan $\sigma\in\Sigma_t$:
(a) for fixed $(t,D)$, $J_\sigma(t,D,b)=bJ^1_\sigma(t,D)+(1-b)J^0_\sigma(t,D)$, and
(b) for fixed $(t,b)$, $J_\sigma(t,D,b)=a_\sigma(t,b)+D\,\beta_\sigma(t,b)$ with
$\beta_\sigma=r_1\P_\sigma(E_{\mathrm{rec}})+h\P_\sigma(E_{\mathrm{res}})$, where $E_{\mathrm{rec}}$ is the
event that some recovery (chosen, after an audit, or forced) occurs in periods $t,\dots,T-1$, and
$E_{\mathrm{res}}$ the event that none occurs and $z_T=1$.
\end{lemma}
\begin{proof}
(a) Under a fixed plan the actions are functions of the observations. Hence the conditional law of the future
given $z_t$ does not depend on $b$. Conditioning on $z_t$ gives the mixture.
(b) The law of the observation process under $\sigma$ does not depend on $D$: increments are independent of
everything (A2), activations and monitor outputs depend only on $z$, and actions are functions of
observations. The initial $D$ enters costs only through the first recovery, which costs $r_0+r_1(D+\Xi)$
with $\Xi$ the increments accumulated so far, and, if no recovery occurs, through the terminal cost
$h(D+\Xi_T)z_T$. All later costs start from $D=0$. Hence the realised cost equals
$D\,[r_1\mathbf{1}_{E_{\mathrm{rec}}}+h\mathbf{1}_{E_{\mathrm{res}}}]$ plus a term that does not depend on
$D$. Taking expectations proves the claim.
\end{proof}

\begin{theorem}[recovery structure, A1 to A5, A6 for (iii)]\label{thm:structure}
(i) $V_t(D,\cdot)$ is concave and piecewise linear on $[0,1]$.
(ii) $\Rc_t(D)$ is a closed interval, possibly empty.
(iii) Under A6, $\Rc_t(D+k)\subseteq\Rc_t(D)$ for every integer $k\ge1$.
(iv) If $1\in\Rc_t(D)$, then $\Rc_t(D)=[\theta_r(t,D),1]$, where $\theta_r(t,D)$ denotes the left endpoint,
and every belief at which auditing is strictly better than recovering lies below $\theta_r(t,D)$. If in
addition A6 holds and $1\in\Rc_t(D+k)$, then $\theta_r(t,D+k)\ge\theta_r(t,D)$.
(v) If $d=1$ and $\alpha=0$, $\Ac_t(D)$ is a closed interval, possibly empty.
\end{theorem}
\begin{proof}
(i) By Lemma~\ref{lem:affine}a, $V_t(D,b)=\min_{\sigma\in\Sigma_t}J_\sigma(t,D,b)$ is a minimum of finitely many
affine functions of $b$. This is the classical result of Smallwood and Sondik (1973).
(ii) After a recovery in period $t$ the carrier is clean and the belief is $0$. The rest of the period has
no activation, and the continuation is $\E_\xi V_{t+1}(\xi,\rho_t)$. Hence
$Q_\Ra(t,D,b)=R(D)+(1-b)c_{fr}+\E_\xi V_{t+1}(\xi,\rho_t)$ is affine in $b$. The function
$G(b)=\min(Q_\Wa,Q_\Aa)(t,D,b)$ is a minimum over the plans whose first action is $\Wa$ or $\Aa$. It is therefore
concave and continuous. Then $\Rc_t(D)=\{b:G(b)-Q_\Ra(t,D,b)\ge0\}$ is a superlevel set of a continuous
concave function on $[0,1]$, which is a closed interval. This is the argument of Lovejoy (1987,
Oper.\ Res.\ 35(4)) for the convexity of stopping regions.
(iii) Fix $(t,b)$ and let $\Sigma^{\neg\Ra}_t$ be the plans whose first action is $\Wa$ or $\Aa$. By
Lemma~\ref{lem:affine}b, for $\sigma\in\Sigma^{\neg\Ra}_t$, $J_\sigma(t,D,b)=a_\sigma+D\beta_\sigma$ with
$\beta_\sigma\le r_1[\P(E_{\mathrm{rec}})+\P(E_{\mathrm{res}})]\le r_1$ under A6, because the two events
are disjoint. Also $Q_\Ra(t,D,b)=r_1D+c(t,b)$ with $c$ independent of $D$. Therefore
\[F(D)=\min(Q_\Wa,Q_\Aa)(t,D,b)-Q_\Ra(t,D,b)=\min_{\sigma\in\Sigma^{\neg\Ra}_t}\big[a_\sigma-c+D(\beta_\sigma-r_1)\big]\]
is a minimum of nonincreasing functions of $D$ and is nonincreasing. If $b\in\Rc_t(D+k)$ then $F(D+k)\ge0$, which gives $F(D)\ge0$
and $b\in\Rc_t(D)$.
(iv) A closed interval in $[0,1]$ that contains $1$ has the form $[\theta,1]$. A belief at which auditing
is strictly better than recovering is not in $\Rc_t(D)$ and is therefore below $\theta_r$. If A6 holds
and $1\in\Rc_t(D+k)$, then by (iii) $[\theta_r(t,D+k),1]\subseteq\Rc_t(D)=[\theta_r(t,D),1]$.
(v) With $d=1$ and $\alpha=0$ an audit reveals $z$: the posterior is $1$ with probability $b$ and $0$ with
probability $1-b$. Hence $Q_\Aa(t,D,b)=c_a+b\,M(1)+(1-b)\,M(0)$ with $M(\beta)=\min(Q_\Wa,Q_\Ra)(t,D,\beta)$,
which is affine in $b$. The function $\min(Q_\Wa,Q_\Ra)$ is concave by the argument of (ii). Hence
$\Ac_t(D)=\{b:\min(Q_\Wa,Q_\Ra)(t,D,b)-Q_\Aa(t,D,b)\ge0\}$ is a closed interval.
\end{proof}

\paragraph{Relation to classical results.} Parts (i) and (ii) are the concavity of POMDP values and the
convexity of stopping sets. The reset structure is that of partially observed machine replacement
(Krishnamurthy 2016), which Hammar and Stadler (DSN 2024) use to prove threshold recovery strategies for
intrusion tolerance. Part (iii) concerns monotonicity in the \emph{observed} derived work. It uses the
linearity of replay cost in $D$, and its direction is the opposite of the conjecture of Definition 3.

\begin{proposition}[the two-threshold structure fails, A1 to A5]\label{prop:counter}
Each instance has $T=2$, $\pi=0.5$, deterministic growth $\xi_t\equiv1$, $\rho_t\equiv0$ and $c_{fr}=0$.
(a) Perfect audits, \PropAParams: at $(t,D)=(0,\PropDzero)$ the optimal first action is $\Wa$ at
$b=\PropBa$, $\Aa$ at $\PropBb$, $\Wa$ at $\PropBc$ and $\Ra$ at $\PropBd$.
(b) A6 holds, \PropBParams: the recovery threshold at $t=0$ is $\PropBThrTwo$ for $D=2$ and
$\PropBThrSix$ for $D=6$. At $b=\PropBBelief$ recovery is optimal for $D=2$ and an audit for $D=6$.
(c) Imperfect audits, \PropCParams: at $(t,D)=(0,6)$ the optimal first action is $\Wa,\Aa,\Wa,\Aa,\Ra$ at
$b=\PropCBeliefs$.
\end{proposition}
\begin{proof}
For each listed belief, the expected cost of every deterministic two-period plan was computed by exhaustive
enumeration ($\PropNpol$ plans in (a), $\PropBNpol$ in (b) and (c)), and the minimum over plans with each
first action agrees with exact backward induction to $10^{-9}$ (Table~\ref{tab:prop2}). The thresholds in
(b) are exact. Every plan's cost is affine in $b$, and its two coefficients follow from its costs at two
interior beliefs. The threshold $\theta_r(0,D)$ is the smallest $b$ at which the lower envelope of the
$\PropBNpol$ plans that recover first meets the lower envelope of the others, located by bisection on
these exact envelopes to $10^{-12}$. The values are \PropBThrTwoExact{} for $D=2$ and \PropBThrSixExact{}
for $D=6$. A narrow audit band lies just below the threshold for $D=2$. In (a) the wait region contains $[0,\PropSegWend]$ and
$[\PropSegWtwostart,\PropSegWtwoend]$ with audits in between, and no thresholds $\theta_a\le\theta_r$ exist.
By Theorem~\ref{thm:structure}(v) the audit region in (a) is an interval, and (c) shows that with
imperfect audits it need not be. In (b) the threshold rises with $D$, against the conjecture that it falls.
\end{proof}

\paragraph{Mechanism of (a).} At $b=\PropBc$ the free monitor output of the next period resolves most of
the uncertainty: an alarm raises the belief to \PropPostHi, where recovery is clearly optimal, and no
alarm lowers it to \PropPostLo, where one audit is bought. Waiting saves the audit cost in the alarm branch
at the price of one more unit of replay. At $b=\PropBb$ an alarm only reaches \PropPostMid, which still
needs an audit, and the audit is bought at once. Ross (1971) observed the same failure of the intuitive
three-region structure in a two-state inspection problem.

\begin{table}[h]
\centering
\caption{Exhaustive enumeration for Proposition~\ref{prop:counter}: minimum expected cost over all
deterministic plans by first action.}\label{tab:prop2}
\begin{tabular}{llrrrrl}
\toprule
instance & belief $b$ & plans & best $W$ & best $R$ & best $A$ & optimal \\
\midrule
(a) & 0.20 & 18{,}396 & 4.034 & 14.000 & 4.800 & $W$ \\
(a) & 0.55 & 18{,}396 & 10.030 & 14.000 & 9.700 & $A$ \\
(a) & 0.82 & 18{,}396 & 13.259 & 14.000 & 13.480 & $W$ \\
(a) & 0.95 & 18{,}396 & 14.632 & 14.000 & 15.300 & $R$ \\
(b), $D=2$ & 0.40 & 63{,}756 & 5.298 & 3.000 & 3.746 & $R$ \\
(b), $D=6$ & 0.40 & 63{,}756 & 7.097 & 7.000 & 5.466 & $A$ \\
(c) & 0.05 & 63{,}756 & 1.346 & 14.000 & 1.596 & $W$ \\
(c) & 0.18 & 63{,}756 & 4.618 & 14.000 & 4.500 & $A$ \\
(c) & 0.33 & 63{,}756 & 7.625 & 14.000 & 7.789 & $W$ \\
(c) & 0.60 & 63{,}756 & 12.930 & 14.000 & 11.460 & $A$ \\
(c) & 0.95 & 63{,}756 & 16.190 & 14.000 & 14.250 & $R$ \\
\bottomrule
\end{tabular}

\end{table}

\paragraph{Search statistics.} Over $\TOneInstances$ random instances ($T\in\{3,\dots,8\}$, $\TOneCells$
cells $(t,D)$, a grid of 1{,}001 beliefs), the solver found no violation of Theorem~\ref{thm:structure}(ii),
no violation of (iii) among $\TOneMonoCells$ comparisons in instances satisfying A6, and no violation of (v)
in $\TOnePAInstances$ random perfect-audit instances. The two-threshold order was broken in
$\PTwoPatternInst$ instances, the audit region was disconnected in $\PTwoAuditInst$, the recovery threshold
increased with $D$ in $\PTwoThrRInst$ and the first non-wait threshold in $\PTwoThrAInst$. Exact recursion
confirmed $\PTwoConfirmed$ of the $\PTwoTried$ flagged cells it re-checked: the first
$\PTwoRecheckPattern$ of the order-break instances and all $\PTwoRecheckAudit$ instances with a
disconnected audit region. The increasing thresholds were not re-checked exactly.

\section{Exact audit timing in the perfect-audit submodel}

Under A7 write $\pi=\pi_0$, $D_s=D_0+\Delta s$, $q_\tau=\pi(1-\lambda)^\tau+(1-\pi)$ (the probability of no
activation before period $\tau$), $N_T=\sum_{s=0}^{T-1}(1-\lambda)^s=(1-(1-\lambda)^T)/\lambda$ and
$g_0=r_1\Delta+\lambda\phi H$.

\begin{lemma}[reduction to first-intervention policies, A1 to A5 and A7]\label{lem:reduce}
There is an optimal policy of the following form: wait until a first intervention time
$\tau\in\{0,\dots,T-1\}$ (if no activation occurred before $\tau$) and then either audit and recover on a
positive outcome or recover blind, or never intervene.
\end{lemma}
\begin{proof}
Since $\sigma_0=\sigma_1$, monitor outputs are independent of the taint and of activations. A policy that
conditions its first intervention on them induces a probability mixture over first-intervention times, and
its expected cost is the same mixture of the costs of deterministic first-intervention times. It is therefore no
better than the best of them. Before the first intervention the only informative event is the absence of
activation. After an activation the mandatory recovery leaves a clean carrier that stays clean
($\rho_t\equiv0$), and the remaining cost is $0$ under $\Wa$. After a blind recovery the same holds. After a
negative audit the carrier is known clean ($d=1$) and the remaining cost is $0$. After a positive audit the
carrier is known tainted ($\alpha=0$). Recovering at once costs $R(D_\tau)$. Any other continuation either
recovers later at cost $R(D_s)\ge R(D_\tau)$, possibly after paying $\phi H$, or never recovers and pays
$hD_T\ge R(D_T)\ge R(D_\tau)$ by A7 and the monotonicity of $R$. Immediate recovery is therefore optimal.
\end{proof}

Under Lemma~\ref{lem:reduce} the candidate costs are
\begin{align*}
J_A(\tau)&=\pi\Big[\sum_{s=0}^{\tau-1}\lambda(1-\lambda)^s\big(\phi H+R(D_{s+1})\big)+(1-\lambda)^\tau R(D_\tau)\Big]+c_aq_\tau,\\
J_B(\tau)&=\pi\sum_{s=0}^{\tau-1}\lambda(1-\lambda)^s\big(\phi H+R(D_{s+1})\big)+q_\tau R(D_\tau),\\
J_N&=\pi\Big[\sum_{s=0}^{T-1}\lambda(1-\lambda)^s\big(\phi H+R(D_{s+1})\big)+(1-\lambda)^ThD_T\Big].
\end{align*}

\begin{theorem}[exact audit timing, A1 to A5 and A7]\label{thm:timing}
(i) For $0\le\tau\le T-2$, $J_A(\tau+1)-J_A(\tau)=\pi(1-\lambda)^\tau g$ with $g=r_1\Delta+\lambda(\phi H-c_a)$.
(ii) An optimal policy is one of four corners: audit at $0$, blind recovery at $0$, audit at $T-1$, never
intervene. If $g>0$ the best audit time is $0$, and if $g<0$ it is $T-1$.
(iii) Let $\pi>0$ and $T\ge2$. Auditing at $0$ is at least as good as every later audit for all
$r_1\ge0$ if and only if $\phi H\ge c_a$. If $\phi H<c_a$, it is strictly better than every later audit exactly when
$r_1>r_1^\star=\lambda(c_a-\phi H)/\Delta$, and the two tie at $r_1=r_1^\star$. As $r_1\to0$,
$g\to\lambda(\phi H-c_a)$.
(iv) The myopic rule of Definition 4 audits at $t=0$ if $\pi g_0\ge c_a$ and never otherwise. If it audits,
$J_{\mathrm{my}}-J^\star=\max\{0,c_a-(1-\pi)R(D_0)\}\le c_a$. If it defers,
\[0\le J_{\mathrm{my}}-J^\star<c_aN_T+\pi(1-\lambda)^T\big(hD_T-R(D_T)\big),\qquad N_T\le\min\{T,1/\lambda\},\]
and the supremum, over all instances satisfying A7 in which the rule defers, of the ratio of the gap to the
right-hand side equals $1$.
\end{theorem}
\begin{proof}
(i) Subtracting, the terms with $s<\tau$ cancel and
\begin{align*}
J_A(\tau+1)-J_A(\tau)&=\pi(1-\lambda)^\tau\big[\lambda(\phi H+R(D_{\tau+1}))+(1-\lambda)R(D_{\tau+1})-R(D_\tau)\big]+c_a(q_{\tau+1}-q_\tau)\\
&=\pi(1-\lambda)^\tau\big[\lambda\phi H+r_1\Delta\big]-c_a\pi\lambda(1-\lambda)^\tau,
\end{align*}
using $R(D_{\tau+1})-R(D_\tau)=r_1\Delta$ and $q_{\tau+1}-q_\tau=-\pi\lambda(1-\lambda)^\tau$.
(ii) By (i) the increments of $J_A$ have the sign of $g$ for every $\tau$, and the minimum over $\tau$ is at
$\tau=0$ if $g>0$ and at $\tau=T-1$ if $g<0$ (any $\tau$ if $g=0$). For blind recovery,
$J_B(\tau+1)-J_B(\tau)=\pi\lambda(1-\lambda)^\tau\phi H+q_\tau r_1\Delta\ge0$. Hence
$\min_\tau J_B(\tau)=J_B(0)=R(D_0)$. With Lemma~\ref{lem:reduce} the optimum is the minimum of $J_A(0)$,
$J_A(T-1)$, $J_B(0)$ and $J_N$.
(iii) With $\pi>0$ and $T\ge2$ there is at least one later audit time and the increments of (i) have the
sign of $g$. Hence auditing at $0$ is at least as good as every later audit iff $g\ge0$, and strictly better
iff $g>0$. For all $r_1\ge0$, $g\ge0$ iff $\lambda(\phi H-c_a)\ge0$. If $\phi H<c_a$, $g>0$ iff
$r_1>r_1^\star$ (using $\Delta>0$), with equality at $r_1^\star$.
(iv) Before any intervention no activation has been observed, and $b_t=\pi(1-\lambda)^t/q_t$ decreases in
$t$. With $d=1$ and $\alpha=0$ the rule audits at $t$ if $b_tg_0\ge c_a$, hence at $0$ if $\pi g_0\ge c_a$ and
never otherwise. After a positive audit $b=1$ and its recovery test $g_0\ge0$ holds. Thus
$J_{\mathrm{my}}\in\{J_A(0),J_N\}$. Telescoping
$\sum_{s<T}\lambda(1-\lambda)^sR(D_{s+1})+(1-\lambda)^TR(D_T)-R(D_0)=r_1\Delta N_T$ and
$\lambda N_T=1-(1-\lambda)^T$ give
\[J_N-J_A(0)=\pi g_0N_T+\pi(1-\lambda)^T\big(hD_T-R(D_T)\big)-c_a .\]
Next, $J_A(0)-J^\star\le c_a$ in every case. First, $J_A(0)-J_B(0)=c_a-(1-\pi)R(D_0)\le c_a$. Second,
$J_A(0)-J_A(T-1)=\pi(\lambda c_a-g_0)N_{T-1}\le\pi c_a\lambda N_{T-1}\le c_a$. Third, $J_A(0)-J_N=c_a-\pi g_0N_T-\pi(1-\lambda)^T(hD_T-R(D_T))\le c_a$, because $g_0\ge0$ and A7 gives
$hD_T\ge R(D_T)$. If the rule defers, $\pi g_0<c_a$ and
$J_{\mathrm{my}}-J^\star=(J_N-J_A(0))+(J_A(0)-J^\star)<c_a(N_T-1)+\pi(1-\lambda)^T(hD_T-R(D_T))+c_a$. If it
audits, $\pi g_0\ge c_a$, and the identity with $N_T\ge1$ and A7 gives $J_N\ge J_A(0)$. Also
$g_0\ge c_a/\pi\ge\lambda c_a$, which gives $J_A(0)\le J_A(T-1)$ and $J^\star=\min\{J_A(0),J_B(0)\}$.
\emph{Supremum.} Take $r_1=0$, $\phi=1$, $D_0=1$, $\pi\lambda H=(1-\epsilon)c_a$, any $r_0>0$ and $h=r_0$.
These choices satisfy A7. The rule defers. Then $J^\star\le J_B(0)=r_0$ and
$J_{\mathrm{my}}=J_N\ge\pi\phi H(1-(1-\lambda)^T)=(1-\epsilon)c_aN_T$. The gap is therefore at least
$(1-\epsilon)c_aN_T-r_0$, while the right-hand side equals $c_aN_T+\pi(1-\lambda)^Tr_0\Delta T$. As $r_0\to0$
the ratio tends to at least $1-\epsilon$, and $\epsilon$ is arbitrary. The ratio never exceeds $1$ by the
bound itself. \emph{Unbounded residual term.} Take $\pi=1$, $r_1=0$, $\phi=1$, $\lambda H=(1-\epsilon)c_a$ and
any $h$ with $r_0\le hD_0$, which satisfies A7. The rule defers, $J^\star\le J_B(0)=r_0$ and
$J_N=(1-\epsilon)c_aN_T+(1-(1-\lambda)^T)r_0+(1-\lambda)^ThD_T$. The gap is at least
$(1-\epsilon)c_aN_T+(1-\lambda)^T(hD_T-R(D_T))$, which grows without bound in $h$, and its ratio to the
right-hand side of (iv) is at least $1-\epsilon$ for every $h$.
\end{proof}

\paragraph{Relation to classical results.} Part (iii) is an analogue of the irreversibility effect of Arrow
and Fisher (1974) and Henry (1974): deferring an audit keeps the option to recover, and the option is
worthless for the irreversible fraction $\phi$ of the harm. No quasi-option-value construction is involved.
The term $-\lambda c_a$ is the continuation value the myopic rule drops: activation itself reveals the taint
for free.

\paragraph{Numerical check.} Over $\TTwoInstances$ random instances satisfying A7 the corner formula matched
exact backward induction with maximal absolute error \TTwoMaxErr{}, the best audit time was always the corner
predicted by the sign of $g$ ($\TTwoCornerViol$ violations), the myopic rule deferred in $\TTwoDefers$
instances, its exact value matched the closed form in $\TTwoMyMatch$ instances, and the bound of (iv) was
never violated ($\TTwoBoundViol$ violations). In the family of the supremum argument with $\epsilon=0.05$ the
ratio of gap to bound is \TTwoTightRatio{} for $r_0=0.01$.

\section{Attacker shift and coverage reallocation}

Under A8 and A9 let $q_i(x)=(1-dx)(1-\lambda_i)$ and $N_i(x)=\sum_{t=0}^{T-1}q_i(x)^t$. A9 makes recovery
after a detection mandatory, as A4 does after an activation. The condition $h\ge r_1$ bounds the price of
this choice: a detected taint left in place is later recovered at $R(D_s)\ge R(D_t)$ (after a detection or
an activation) or reaches $T$ and costs $hD_T\ge r_1D_T\ge r_1D_t=R(D_t)-r_0$. Leaving it therefore never
saves more than the fixed price $r_0$.

\begin{theorem}[attacker shift and reallocation, A8, A9, and A10 for (iii)]\label{thm:game}
(i) $L_i(x,r_1)=F_i(x)+r_1E_i(x)$ with
\begin{align*}
F_i(x)&=(1-dx)\lambda_i\phi_iH\,N_i(x)+r_0\big(1-q_i(x)^T\big)+h\Delta_iT\,q_i(x)^T,\\
E_i(x)&=\textstyle\sum_{t=0}^{T-1}q_i(x)^t\big[dx\,\Delta_it+(1-dx)\lambda_i\Delta_i(t+1)\big].
\end{align*}
Both are nonnegative polynomials in $x$. At $x=0$ the irreversible term of $F_i$ equals
$\phi_iH(1-(1-\lambda_i)^T)=\lambda_i\phi_iH\,T\,(1+O(\lambda_iT))$.
(ii) For any fixed commitment $x$, if carrier $i$ is a best response at $r_1$ and $j$ at $r_1'>r_1$
(abstention counts as a carrier with $sE=0$), then $s_jE_j(x_j)\ge s_iE_i(x_i)$. At $r_1=0$ the best
response maximises $s_iF_i(x_i)-\kappa_i$ if that maximum is positive, and abstains otherwise.
(iii) Under A10 the minimax commitment is unique and differentiable at $r_1$, and
\[\frac{dx_i^*}{dr_1}=\frac{s_iE_i(x_i^*)-\bar E}{|\partial U_i/\partial x_i|}\ (i\in S),\qquad
\frac{du^*}{dr_1}=\bar E=\frac{\sum_{j\in S}w_js_jE_j(x_j^*)}{\sum_{j\in S}w_j}\ge0,\]
with $w_j=1/|\partial U_j/\partial x_j|$, and $dx_j^*/dr_1=0$ for $j\notin S$.
\end{theorem}
\begin{proof}
(i) Condition on the first event. In period $t$ the carrier is still tainted with probability $q^t$. It is
then detected with probability $dx$ (cost $R(\Delta_it)$), or else activates with probability $\lambda_i$
(cost $\phi_iH+R(\Delta_i(t+1))$). With probability $q^T$ it survives to $T$ (cost $h\Delta_iT$). Summing and
separating the coefficient of $r_1$ gives $L_i=F_i+r_1E_i$. The $r_0$ terms sum to $r_0\sum_tq^t(1-q)=r_0(1-q^T)$
because $dx+(1-dx)\lambda_i=1-q$. At $x=0$, $\lambda_iN_i(0)=1-(1-\lambda_i)^T$, and the expansion follows from
the binomial series.
(ii) For fixed $x$, $U_k(r_1)=[s_kF_k(x_k)-\kappa_k]+r_1s_kE_k(x_k)$ is affine in $r_1$ with slope
$\beta_k=s_kE_k(x_k)$. From $U_i(r_1)\ge U_j(r_1)$ and $U_j(r_1')\ge U_i(r_1')$, adding gives
$(\beta_j-\beta_i)(r_1'-r_1)\ge0$.
(iii) \emph{Smoothness.} $U_i(x,r_1)=s_i(F_i(x)+r_1E_i(x))-\kappa_i$ is a polynomial in $x$ and affine in $r_1$,
hence jointly $C^\infty$. Every feasible coverage lies in $[0,\min\{1,B\}]$, where A10 and compactness give
$\partial_xU_i\le-\delta<0$. For $r_1'$ near $r_1$ the inequality $\partial_xU_i(\cdot,r_1')<0$ persists
because $\partial_xE_i$ is bounded. \emph{Water filling.} Because $u^*>0$, the objective $\max\{0,\max_iU_i\}$ equals $\max_iU_i$ near
$x^*$, and lowering the maximum payoff lowers the objective. At a minimax commitment every carrier with
positive coverage attains the maximum. Otherwise move a small amount $\varepsilon$ of its coverage, in equal parts, to
the carriers that attain the maximum. Their coverage stays in $[0,\min\{1,B\}]$ for small $\varepsilon$ (it is
below $1$ under A10 and the total is unchanged), and each of their payoffs falls strictly while the donor stays
below the maximum, which contradicts optimality. In particular $x_j^*=0$ for $j\notin S$. \emph{Uniqueness.}
Let $x'$ be any minimax commitment. For $j\notin S$, $U_j(x'_j)\le U_j(0)<u^*$. For $i\in S$,
$U_i(x'_i)\le u^*=U_i(x^*_i)$ gives $x'_i\ge x^*_i$, and $\sum_{i\in S}x'_i\le B=\sum_{i\in S}x^*_i$ forces
$x'=x^*$. This is the equalisation of ORIGAMI (Kiekintveld et al.\ 2009). \emph{Local constancy of $S$.} The
value $v(r_1)=\min_{x\in X_B}\max_iU_i(x_i,r_1)$ is continuous, it stays positive near $r_1$, and the solution correspondence is upper
hemicontinuous (Berge's maximum theorem, continuous objective, fixed compact feasible set). For $r_1'$ near
$r_1$ every minimax commitment $x(r_1')$ is therefore close to $x^*$. For $j\notin S$, $U_j(0,r_1')<v(r_1')$ by
continuity. Then $j$ does not attain the maximum and water filling gives $x_j(r_1')=0$. For $i\in S$,
$x_i(r_1')$ is close to $x^*_i\in(0,1)$ and positive, and $i$ attains the maximum. The budget binds at
$r_1'$, since otherwise raising the coverage of the maximisers would lower the maximum. \emph{Implicit
function.} Near $(x^*_S,u^*,r_1)$ the system $U_i(x_i,r_1)-u=0$ for $i\in S$ and $\sum_{i\in S}x_i-B=0$ has a
Jacobian with respect to $(x_S,u)$ with determinant $\pm\prod_{i\in S}\partial_xU_i\cdot\sum_{i\in S}1/\partial_xU_i\neq0$.
The implicit function theorem gives a unique $C^1$ solution near $r_1$, which by local constancy and
uniqueness is the minimax commitment. Differentiating $U_i(x_i^*(r_1),r_1)=u^*(r_1)$ gives
$\partial_xU_i\,x_i^{*\prime}+s_iE_i=u^{*\prime}$. Summing $x_i^{*\prime}=(s_iE_i-u^{*\prime})/|\partial_xU_i|$
over $S$ and using $\sum_ix_i^{*\prime}=0$ yields $u^{*\prime}=\bar E\ge0$.
\end{proof}

\paragraph{Relation to the classical result.} ORIGAMI and the analysis of Korzhyk et al.\ (2011) equalise
attacker utilities over the attack set for fixed target values. Here the values depend on the recovery
infrastructure, and part (iii) gives the first-order shift of the minimax coverage when the price of rollback
changes. Its sign is set by the reversible exposure $s_iE_i$ relative to the weighted mean $\bar E$. The
irreversible fraction $\phi_i$ enters only through $F_i$, which sets $x^*$ and with it $E_i(x_i^*)$, the
weights and $\bar E$. It can change the sign only through these quantities. The same structure appears in
inspection games (Avenhaus, von Stengel and Zamir 2002) with fixed values.

\begin{proposition}[no simple migration to $\arg\max\phi$, A8 and A9]\label{prop:p3}
(a) In the four-carrier instance of the paper ($h=10$) at $r_1=0$ without coverage, the best response is the
memory note with $\phi=0$ (payoff $\PThreeUmemory$), and the queued call with $\phi=1$ yields the least
($\PThreeUqueue$). (b) The carrier with the largest $\phi$ can lose coverage when rollback gets cheaper.
(c) Where the attack set changes, the minimax coverage need not be differentiable.
\end{proposition}
All payoffs of (a) use $T=20$, $H=20$, $d=0.9$ and $r_0=3$. With $h=1$, the residual price of E1 and E3,
the best response without coverage is the \PThreeLowHbr{} (payoff \PThreeLowHbrU) and the
\PThreeLowHleast{} still yields the least (\PThreeLowHleastU). The point of (a), that the attacker does not
go to the largest $\phi$, therefore holds at both residual prices.
\begin{proof}
(a) is read from the closed forms of Theorem~\ref{thm:game}(i). (b) Take carriers
$a=(\Delta,\lambda,\phi,s,\kappa)=(1,0.02,1,0.5,1)$ and $b=(0.1,0.2,0,0.8,0.5)$ with $T=20$, $H=20$, $h=2$,
$r_0=1$, $d=0.9$ and $B=0.5$, for which A9 holds when $r_1\in[0,2]$. On a grid of 21 values of $r_1\in[0,2]$ both
carriers stay in the attack set and A10's derivative condition holds. The coverage of $a$ falls from
$\PThreeExXaTwo$ at $r_1=2$ to $\PThreeExXaZero$ at $r_1=0$. Its reversible exposure $s_aE_a$ is
$\PThreeExSBaZero$ against $\PThreeExSBbZero$ for $b$ at $r_1=0$, a ratio of $\PThreeExRatio$, and
Theorem~\ref{thm:game}(iii) moves coverage away from it as $r_1$ falls. In addition, among the
$\PThreeAdm$ random games that satisfy A9 at $r_1=2$ ($h\ge2$) and have a positive attacker level at
$r_1=0$ and $r_1=2$, the largest-$\phi$ carrier received less coverage at $r_1=0$ than at $r_1=2$ in
$\PThreeAdmLoses$. (c) In the four-carrier instance the
queued call leaves the attack set at $r_1\approx\PThreeKinkR$ (bisection), where the one-sided
finite-difference derivatives of its coverage are \PThreeKinkLeft{} and \PThreeKinkRight.
\end{proof}

\paragraph{Numerical check.} The closed form of (i) agreed with Monte Carlo within $3.5$ standard errors in
$\TThreeMCok$ of $\TThreeMCn$ cells. Best-response slopes were monotone in $r_1$ in all $\TThreeBRn$ random
games ($\TThreeBRviol$ violations), a check of the implementation of an algebraic identity. Under
A10, checked on a 401-point coverage grid, the derivative formula matched central finite differences within
$5\%$ in $\TThreeDok$ of $\TThreeDn$ games with matching signs in $\TThreeDsign$. $\TThreeAfail$ random games
failed the other conditions of A10 and $\TThreeNonpos$ more had a nonpositive level, and all were excluded.
In the four-carrier instance A10 holds at every grid value of $r_1$ (the queued call enters the attack set
between grid values, near $r_1\approx\PThreeKinkR$).

\section{Accuracy of the belief grid}

The planner tables use linear interpolation on a uniform belief grid. Against the exact recursion over the
reachable belief tree, on the E1 base instance with $T\le\GridMaxT$, every derived work $D\le D_0$ and 37
beliefs (most of them off the grid), the relative error of the 1{,}001-point grid is below \GridMaxRel.
Longer horizons are out of reach of the exact recursion.

At the horizon of E3 ($T=24$) we compare every ROVER table (four carriers, $r_1\in\{0,0.1,1,10\}$) with an
\PostGridRef-point grid over all periods, every derived work $D\le t$ and 1{,}001 beliefs. The 501-point
tables differ in value by at most \PostGridMaxErr{} (\PostGridMaxRel\% relative) and in the chosen action in at
most \PostGridActDis\% of these states, and the errors fall as the grid is refined. With \PostGridDense-point
tables ROVER takes the same decisions, with the same losses, in all \PostGridIdent{} held-out main episodes at
the four replay prices. This is a convergence study against a finer grid and not an error bound.

\section{Sensitivity grid of E1}

Table~\ref{tab:esens} varies one parameter of the E1 base instance at a time and reports the optimal audits
per 100 periods. The pattern of the main text (with $\phi=1$ audits are worth more than with $\phi=0$ when
rollback is free, and with $\phi=0$ the audit rate rises with $r_1$) holds in \EOneSensHolds{} of the
\EOneSensTotal{} distinct configurations. It breaks for \EOneSensBreaks. With dear or weak audits the
$\phi=1$ policy replaces audits by precautionary recoveries when rollback is free, and with rare
activations neither policy intervenes at $r_1=10$, where a recovery costs more than the expected harm.

\begin{table}[h]
\centering
\caption{E1 sensitivity grid: optimal audits per 100 periods for $\phi\in\{0,1\}$ and $r_1\in\{0,1,10\}$,
one parameter changed at a time from the base instance.}\label{tab:esens}
\begin{tabular}{lrrrrrrl}
\toprule
& \multicolumn{3}{c}{$\phi=0$, $r_1=0,1,10$} & \multicolumn{3}{c}{$\phi=1$, $r_1=0,1,10$} & pattern \\
\midrule
$c_a=0.25$ & 4.5 & 5.0 & 17.6 & 17.9 & 14.5 & 23.6 & holds \\
$c_a=0.5$ & 3.5 & 2.7 & 8.9 & 15.7 & 11.8 & 14.0 & holds \\
$c_a=1$ & 1.6 & 1.9 & 8.1 & 5.3 & 7.7 & 9.5 & holds \\
$c_a=2$ & 0.0 & 0.0 & 5.6 & 0.0 & 3.6 & 7.1 & breaks \\
$d=0.5$ & 0.0 & 0.0 & 7.2 & 0.0 & 3.8 & 8.1 & breaks \\
$d=0.7$ & 0.7 & 0.8 & 7.3 & 0.4 & 5.3 & 8.8 & breaks \\
$d=0.9$ & 1.6 & 1.9 & 8.0 & 5.3 & 7.7 & 9.5 & holds \\
$d=1$ & 2.0 & 1.9 & 8.0 & 6.7 & 7.6 & 8.7 & holds \\
$\lambda=0.02$ & 2.0 & 3.1 & 0.0 & 2.2 & 4.1 & 0.0 & breaks \\
$\lambda=0.05$ & 1.9 & 2.5 & 5.7 & 2.7 & 4.6 & 6.9 & holds \\
$\lambda=0.1$ & 1.6 & 1.9 & 8.0 & 5.3 & 7.7 & 9.4 & holds \\
$\lambda=0.2$ & 1.1 & 0.7 & 9.1 & 11.5 & 8.1 & 11.1 & holds \\
$\sigma_1=0.1$ & 4.0 & 0.0 & 10.8 & 13.4 & 10.5 & 11.2 & holds \\
$\sigma_1=0.25$ & 3.2 & 1.4 & 8.0 & 12.0 & 8.9 & 11.8 & holds \\
$\sigma_1=0.4$ & 1.6 & 1.9 & 8.1 & 5.3 & 7.7 & 9.5 & holds \\
$\sigma_1=0.7$ & 0.7 & 0.3 & 6.0 & 1.4 & 4.1 & 6.8 & holds \\
\bottomrule
\end{tabular}

\end{table}

\section{Development history and mechanics of E3}

On the development streams, our first full run used a constant insertion hazard and a replay-only
executor. It lost to rollback only at $r_1=0$ (\EThreeVOneRoverZero{} against \EThreeVOneRollbackZero) and
to reset all at $r_1=10$ (\EThreeVOneRoverTen{} against \EThreeVOneResetTen). Both changes are kept as
ablations, and on held-out data each original loses at its extreme: without the window at $r_1=0$
(\EThreeLZeroConst{} against \EThreeLZeroRover, Holm $p\EThreePZeroConst$) and with replay only at $r_1=10$
(\EThreeLTenReplayOnly{} against \EThreeLTenRover).

The held-out run issued \EThreeGitCalls{} git commands. A recovery took \EThreeMsPerRecovery\,ms of git time
on average and re-executed or cherry-picked the work since the checkpoint (\EThreeRegen{} re-executions,
\EThreeCherry{} cherry-picks, \EThreeConflicts{} conflicts). Irreversible pushes delivered
\EThreeRemoteCanary{} canary lines to the bare remote, and no policy could remove them.

\section{Exploratory analyses added after the confirmatory run}

The analyses below were run after the confirmatory E3 run and are not part of its protocol. Their
definitions are our choices and are stated here in full.

\paragraph{One-step controllers in E1.} The myopic rule (Definition~4) waits or audits and never recovers
blindly. A one-step controller chooses $\arg\min\{Q_W,Q_R,Q_A\}$ with the one-period transitions of
Definition~1 and a fixed continuation value $V_{t+1}$ in place of the optimal cost-to-go,
and after an audit it recovers when its own $Q_R\le Q_W$ at the posterior. With the terminal cost
$hDb$ as continuation it is a plain one-step lookahead. With the exact cost-to-go of never auditing it is
one policy-improvement step on that policy (a rollout). Each is run with and without blind recovery and
evaluated exactly on the belief grid as the myopic rule is, with a Monte Carlo cross-check
(largest $|z|=\PostOsMcZ$). Table~\ref{tab:revos} gives the gaps to the optimum on the E1 grid. Allowing
blind recovery helps the rollout controller in \PostOsBlindHelps{} and hurts it in \PostOsBlindHurts{} of the
\PostOsCells{} configurations of the E1 grid and the sensitivity grid, and the rollout controller beats the
myopic rule in \PostOsRollBeats{} of them.

\begin{table}[h]
\centering
\caption{E1, gap to the optimal loss in \% for the myopic rule, the myopic rule with the DP's recovery
decisions, the one-step lookahead with blind recovery, and the rollout controller without and with blind
recovery.}\label{tab:revos}
\begin{tabular}{rrrrrrrr}
\toprule
$\phi$ & $r_1$ & DP loss & Myopic & Myopic, DP rec. & One-step & Rollout W/A & Rollout W/R/A \\
\midrule
0 & 0 & 2.79 & 69.6 & 69.6 & 61.8 & 40.4 & 42.5 \\
0 & 0.1 & 3.60 & 44.9 & 44.9 & 35.0 & 25.5 & 24.8 \\
0 & 1 & 9.15 & 5.0 & 5.0 & 5.0 & 0.8 & 0.8 \\
0 & 10 & 47.00 & 17.7 & 5.1 & 13.7 & 0.7 & 0.7 \\
0.5 & 0 & 5.36 & 56.1 & 56.1 & 10.2 & 22.2 & 25.4 \\
0.5 & 0.1 & 6.08 & 45.5 & 45.5 & 6.4 & 17.4 & 17.4 \\
0.5 & 1 & 11.12 & 6.5 & 6.4 & 19.0 & 4.9 & 4.9 \\
0.5 & 10 & 48.73 & 18.0 & 6.2 & 17.1 & 0.9 & 0.9 \\
1 & 0 & 7.05 & 24.2 & 23.9 & 2.4 & 7.1 & 10.5 \\
1 & 0.1 & 7.70 & 19.7 & 19.7 & 1.6 & 5.8 & 5.8 \\
1 & 1 & 12.68 & 7.2 & 7.2 & 28.3 & 2.8 & 2.8 \\
1 & 10 & 50.38 & 17.7 & 6.5 & 20.4 & 1.0 & 1.0 \\
\bottomrule
\end{tabular}

\end{table}

\paragraph{Publication harm in E3.} The main runs have no benign releases. We add releases at rate
$\nu\in\{1/24,1/4\}$ per milestone, with no gate, to the main condition at $r_1=1$ for nine policies on
the held-out cells (\PostPubEpisodes{} episodes). A release publishes the work branch to the bare remote.
Every decision count and loss equals the corresponding main episode, and the leaks can therefore be
priced after the fact. We count the distinct canary lines that releases publish and that no irreversible activation
published before, and charge either $h_{\mathrm{pub}}\in\{1,2,5,20\}$ per such line or
$H_{\mathrm{rel}}\in\{1,5,20\}$ per release that contains one. Table~\ref{tab:revpub} lists the leaks and
the losses at the highest prices. With rare releases ROVER has the lowest loss at every price. With
frequent releases it has the lowest loss up to a price of five per line or per release and drops to
rank \PostPubFreqWorstRank{} at a price of 20, behind the immutable validator and periodic review. The testbed's
counter of canary occurrences over remote refs exceeds the number of distinct leaked lines up to
\PostPubRatioMax-fold (\PostPubRatioRo-fold for ROVER) with frequent releases.

\begin{table}[h]
\centering
\caption{E3 with releases: distinct newly leaked canary lines per episode, share of episodes with a leak,
and mean loss with frequent releases without a publication price, at 20 per leaked line and at 20 per
contaminated release.}\label{tab:revpub}
\begin{tabular}{lrrrrrrr}
\toprule
 & \multicolumn{2}{c}{$\nu=1/24$} & \multicolumn{2}{c}{$\nu=1/4$} & \multicolumn{3}{c}{loss, $\nu=1/4$} \\
\cmidrule(lr){2-3}\cmidrule(lr){4-5}\cmidrule(lr){6-8}
Policy & lines & leak \% & lines & leak \% & none & 20/line & 20/release \\
\midrule
ROVER & 0.26 & 8 & 1.02 & 30 & 9.28 & 29.70 & 24.61 \\
ROVER, no audits & 0.55 & 15 & 1.58 & 39 & 10.37 & 41.95 & 33.37 \\
Rollback only & 0.33 & 12 & 1.22 & 38 & 11.36 & 35.77 & 30.61 \\
No audit & 0.70 & 20 & 2.00 & 46 & 11.86 & 51.86 & 45.20 \\
Random audit & 0.69 & 19 & 1.85 & 45 & 13.46 & 50.46 & 45.29 \\
Myopic rule & 0.67 & 18 & 1.83 & 45 & 11.50 & 48.08 & 40.66 \\
Periodic review & 0.06 & 4 & 0.34 & 22 & 16.07 & 22.82 & 23.07 \\
Immutable validator & 0.02 & 2 & 0.04 & 5 & 15.19 & 15.94 & 16.86 \\
Reset all & 0.30 & 10 & 0.98 & 34 & 25.94 & 45.60 & 40.77 \\
\bottomrule
\end{tabular}

\end{table}

\paragraph{Cost semantics in E3.} An instrumented re-run of ROVER and tuned rollback only at
$r_1\in\{0.1,1\}$ (\PostCsEpisodes{} episodes, decisions equal to the main runs) logs the planner's price
$r_0+\min(r_1,w)D$ at every recovery and the executed price. Executed recoveries cost at most
\PostCsCostErrMax\% more on average, because cherry-pick conflicts force \PostCsConfMin{} to
\PostCsConfMax{} extra re-executions per recovery. For carriers still tainted at the end, the surviving canary
lines fall \PostCsResLowMin{} to \PostCsResLowMax\% short of the terminal prediction $hP_{\mathrm{infl}}D$,
because $D$ counts dependent commits made before the insertion. Measured against $hP_{\mathrm{infl}}$
times the direct reads after the insertion, the ratio is \PostCsDirectMin{} to \PostCsDirectMax. Transitive dependencies make $D$ grow
\PostCsGrowMin{} to \PostCsGrowMax\% faster than the planner's growth model for the carriers that are not read
at every milestone.

\paragraph{Estimands.} Next to the signed-rank test we test the mean of the paired per-stream differences
with a paired $t$-test and a sign-flip permutation test (200{,}000 sign vectors). The primary decisions do
not change (Holm $p\PostStHOneT$ and $p\PostStHOneSf$ for H1, $p\PostStHTwoT$ and $p\PostStHTwoSf$ for H2).
Over the \PostStTests{} comparisons against ROVER, \PostStFlips{} Holm decision changes. $A_{12}$ as reported
compares every stream of one policy with every stream of the other (\PostStHOneAtw{} for H1), while the
paired share of streams on which ROVER is lower is \PostStHOneWin.

\paragraph{Adaptive attacker.} The attacker's payoff per carrier rests on \PostAdNMin{} to \PostAdNMax{}
attacked episodes, and the reported maximum over four carriers is biased upward. Choosing the carrier on
seeds 8 to 11 and scoring it on seeds 12 to 15, and the reverse, gives \PostAdRoSplit{} against ROVER and
\PostAdVaSplit{} against the immutable validator.

\paragraph{Exact check of A10.} For the E2 game and the instance of Proposition~\ref{prop:p3}, $F_i$ and $E_i$ are
polynomials in $x_i$ (degree 20). We build them with exact rational coefficients from the stated
parameters and count the real roots of $\partial U_i/\partial x_i=s_i(F_i'+r_1E_i')$ on $[0,\min\{1,B\}]$ by
Sturm sequences. There is none, and the derivative is negative, at both ends of the range of $r_1$ for every
carrier. The derivative is affine in $r_1$, and the monotonicity part of A10 therefore holds for every $r_1$ in
$[0,10]$ (E2) and $[0,2]$ (Proposition~\ref{prop:p3}). The other parts of A10 concern the minimax solution and remain
numerical checks.

\section*{Use of generative AI}

Generative AI tools were used to assist with drafting and editing parts of the text, writing and
debugging the simulation, solver and testbed code, and checking bibliographic metadata. The authors
designed the study, verified every formal statement and proof, re-ran the artifact that produces all
reported numbers, checked every reference against its publisher or archive record, and take full
responsibility for the content.

\end{document}
